%% -*- coding:utf-8 -*-
\chapter{Verbstellung: Verberst- und Verbzweitstellung}
\label{chap-verb-position}

Dieses Kapitel behandelt die Analyse der Stellung des finiten Verbs in V2-Sprachen. Ich konzentriere mich auf
das \ili{Dänisch}e und das \ili{Deutsch}e, die als prototypische Beispiele dienen können: Das \ili{Dänisch}e ist eine SVO-Sprache, während das \ili{Deutsch}e
SOV ist. Ich diskutiere zunächst Argumente für die Klassifikation des \ili{Deutsch}en als SOV-Sprache und
stelle die nötigen Daten zum \ili{Dänisch}en bereit, um danach die jeweiligen Analysen zu erläutern.

\section{Die Phänomene}

Abschnitt~\ref{sec-intro-svo}\is{Abfolge!SOV|(}\is{Abfolge!SVO|(} enthält eine Diskussion der Grundabfolge von Subjekt, Objekt und Verb in den Sprachen
der Welt und insbesondere in den \ili{germanisch}en Sprachen. Ich habe die Klassifikation diskutiert, die
der World Atlas of Language Structures \citep{Dryer2013a} bereitstellt und die nahelegt, dass das \ili{Deutsch}e eine Sprache ohne dominante
Konstituentenabfolge ist, aber zwei Abfolgen aufweist, die häufig beobachtet werden können: SOV (Nebensätze und
Sätze mit einem Auxiliar) und SVO (in Hauptsätzen ohne Auxiliar).
% Textbooks used in
% \ili{German} schools claim that \ili{German} is a SVO language.\footnote{
% ``Das Subjekt steht normalerweise am Anfang des Satzes. [\ldots] Das Prädikat steht fast immer an
% zweiter Stelle des Satzes. Es kann aus einem Wort oder aus mehreren Wörtern bestehen:
% \emph{kränkelt} -- \emph{langweilt sich} -- \emph{kann spielen}.'' Textbook for the 5th class by
% Westermann Verlag, 2018.
% } 
Im folgenden Unterabschnitt diskutiere ich diese
Annahme im Detail. Abschnitt~\ref{sec-German-as-SOV} erklärt, warum Forscherinnen und Forscher in der
Mainstream Generative Grammar\footnote{
  Der Begriff \emph{Mainstream Generative Grammar} wird für Arbeiten in der Tradition von
  \citegen{Chomsky81a} Lectures on Government \& Binding und \citegen{Chomsky95a-u} Minimalist Program verwendet.
} und auch die meisten Syntaktikerinnen und Syntaktiker, die in anderen Rahmentheorien arbeiten, davon ausgehen,
dass das \ili{Deutsch}e eine SOV-Sprache ist. Abschnitt~\ref{sec-Germanic-SVO-verb-position} behandelt das \ili{Dänisch}e als
einen Vertreter der \ili{germanisch}en SVO-Sprachen und erklärt, wie verb"=initiale Sätze in diesen Sprachen
am besten beschrieben werden können. Abschnitt~\ref{sce-verb-second} ist dem Vorfeld-Voranstellen im \ili{Englisch}en und
den Verbzweitsätzen im \ili{Deutsch}en (und den anderen \ili{germanisch}en Sprachen) gewidmet.

%\subsection{SVO or SOV?}

\subsection{Deutsch als SVO-Sprache?}

Zu behaupten, dass SVO eine Grundabfolge ist, ist auf der Basis reinen Zählens einigermaßen merkwürdig angesichts der
Tatsache, dass die meisten \ili{deutsch}en Sätze das Subjekt ohnehin nicht an erster Stelle haben. Der folgende Text
mag als Beispiel dienen:
\eanoraggedright
\gruenbf{Für selbstfahrende Autos} soll es in Deutschland nach Angaben von Bundes\-ver\-kehrs\-mi\-nis\-ter
Alexander Dobrindt (CSU) bald eine Teststrecke geben. \gruenbf{Auf der Autobahn A9 in Bayern} \rotit{sei
ein Pilotprojekt „Digitales Testfeld Autobahn“ geplant}, wie aus einem Papier des
Bundesverkehrsministeriums hervorgeht. \gruenbf{Mit den ersten Maßnahmen für diese Teststrecke}
solle schon in diesem Jahr begonnen werden. \gruenbf{Mit dem Projekt} soll die Effizienz von
Autobahnen generell gesteigert werden. \gruenbf{„\rotit{Die Teststrecke} soll so digitalisiert und
  technisch ausgerüstet werden, dass es dort zusätzliche Angebote der Kommunikation zwischen Straße
  und Fahrzeug wie auch von Fahrzeug zu Fahrzeug geben wird“}, sagte Dobrindt zur Frankfurter
Allgemeinen Zeitung.  \gruenbf{Auf der A9} sollten sowohl Autos mit Assistenzsystemen als auch
später vollautomatisierte Fahrzeuge fahren können. \gruenbf{Dort} soll die Kommunikation nicht nur
zwischen Testfahrzeugen, sondern auch zwi\-schen Sensoren an der Straße und den Autos möglich sein,
etwa zur Übermittlung von Daten zur Verkehrslage oder zum Wetter. \gruenbf{\rotit{Das Vorhaben}
  solle im Verkehrsministerium von einem runden Tisch mit Forschern und Industrievertretern
  begleitet werden,} sagte Dobrindt. \rotit{Dieser} solle sich unter anderem auch mit den
komplizierten Haftungsfragen beschäftigen.  Also: \rotit{Wer} zahlt eigentlich, wenn ein
automatisiertes Auto einen Unfall baut?  \gruenbf{[\gruenbf{Mithilfe der Teststrecke}] solle die
  deutsche Automobilindustrie auch beim digitalen Auto „Weltspitze sein können“,} sagte der
CSU-Minister. \rotit{Die deut\-schen Hersteller} sollten die Entwicklung nicht Konzernen wie etwa
Google überlassen.  \gruenbf{Derzeit} ist Deutschland noch an das „Wiener Über\-ein\-kom\-men für den
Straßenverkehr“ gebunden, das Autofahren ohne Fahrer nicht zulässt. \gruenbf{Nur unter besonderen
  Auflagen} sind Tests möglich.  \rotit{Die Grünen} halten die Pläne für
unnütz. \rotit{Grünen-Verkehrsexpertin Valerie Wilms} sagte der Saarbrücker Zeitung: „\rotit{Der
  Minister} hat wichtigere Dinge zu erledigen, als sich mit selbstfahrenden Autos zu beschäftigen.“
\rotit{Die Technologie} sei im Verkehrsbereich nicht vordringlich, \gruenbf{auch} stehe sie noch ganz am
Anfang.  \gruenbf{Aus dem grün-rot regierten Baden-Württemberg – mit dem Konzernsitz von Daimler –}
kamen hingegen andere Töne. \gruenbf{\rotit{Was in Bayern funktioniere,} müsse auch in
  Baden-Württemberg möglich sein,} sagte Wirtschaftsminister Nils Schmid (SPD). \gruenbf{Von den
  topografischen Gegebenheiten} biete sich die Autobahn A81 an.\footnote{\emph{Selbstfahrende Autos: A9 soll Teststrecke werden}, taz, 2015-01-27}
\z
%\largerpage
Die Subjekte sind rot markiert und die Nicht-Subjekte grün. Ich habe auch Subjekte/Nicht-Subjekte
innerhalb von eingebetteten Sätzen gezählt. Das Verhältnis liegt bei 11 Subjekten (einschließlich eines Subjektsatzes) gegenüber 16
Nicht-Subjekten (\emph{ein automatisiertes Auto} und \emph{Autofahren ohne Fahrer} in den SOV-Sätzen wurden nicht mitgezählt, aber diese sind
natürlich ebenfalls Gegenbeispiele zur SVO-Behauptung). Die Frage ist also: Was sagt uns
diese Zahl? Natürlich könnten wir die grammatischen Funktionen des
vorangestellten Materials nun weiter differenzieren. Wir würden feststellen, dass 3 Objektsätze vorangestellt sind; der Rest der vorangestellten
Konstituenten sind Adverbiale. Wir könnten schlussfolgern, dass SVO häufiger ist als OVS, aber zu sagen, dass SVO
basal ist, wäre nicht angemessen. Vielmehr sollte AdvVSO als Grundmuster angesehen werden, wenn wir
diesen kleinen Text als unsere empirische Basis annehmen. Natürlich reicht es nicht aus, diesen Text als Grundlage
wissenschaftlicher Behauptungen anzunehmen.
% \citet[\page 162]{Hoberg81a} counted elements in the \vf of
% 4717 sentences from different text types and found that 63.15\,\% of the \vf constituents were
% subjects. However, \citet[\page 163]{Hoberg81a} also pointed out that if one looks at subjects and
% where they appear it is only about half of the subjects that appear in the \vf, the other half are
% placed in the \mf. See also \citew[\page 1584]{Hoberg97a}.
%\largerpage 
\citet[Abschnitt~4]{HK2005a} hat die \vf-Konstituenten
in den Korpora TüBa-D/S und Z untersucht. Das S-Korpus enthält gesprochenes \ili{Deutsch} aus dem maschinellen
Übersetzungsprojekt \verbmobil und das Z-Korpus Sätze aus der \ili{deutsch}en Zeitung taz. Beide
Korpora sind nach grammatischer Funktion annotiert. Die TüBa-D/S bestand aus insgesamt 38.342 Bäumen
und die Baumbank TüBa-D/Z hatte 22.087 Bäume, als der Aufsatz 2005 geschrieben wurde. Die beiden Korpora hatten
in 50,3\,\% bzw.\ 52,1\,\% der Sätze mit einem \vf ein Subjekt im \vf. 
% Unfortunately, the number
% of subjects in the \mf was not reported. But assuming effects similar to the ones Hoberg observed,
% the number of subjects in the \mf may be larger than the number of fronted subjects. 
SVO als Grundabfolge anzunehmen wäre also
nicht hilfreich, denn in etwa 50\,\% der Sätze und vielleicht sogar noch
mehr hätte man es mit einer Abfolge zu tun, in der das Subjekt nicht an erster Stelle steht. Hinzu
käme das Problem der Nebensätze, die eindeutig eine SOV-Abfolge
haben.\footnote{
Die Sätze in (i) stehen in SOV-Abfolge. (id) ist infinit und hat kein Subjekt.
\eal
\ex
dass es dort zusätzliche Angebote der Kommunikation zwischen Straße und Fahrzeug wie auch von Fahrzeug zu Fahrzeug geben wird
\ex
wenn ein automatisiertes Auto einen Unfall baut
\ex
das Autofahren ohne Fahrer nicht zulässt
\ex
als sich mit selbstfahrenden Autos zu beschäftigen
\zllast
}
 Deshalb gehen Syntaktikerinnen und Syntaktiker verschiedener Rahmentheorien (siehe \citealp{MuellerGT-Eng} für
Ansätze in GB, Minimalismus, LFG, Kategorialgrammatik und HPSG) davon aus, dass SOV die Grundabfolge des
\ili{Deutsch}en ist. Das finite Verb wird vorangestellt, um den Satztyp zu markieren, und eine Konstituente wird vor
dieses Verb gestellt. Diese vorangestellte Konstituente kann das Subjekt, ein Objekt oder eine beliebige andere Konstituente des
Satzes sein. Es kann sogar ein Abhängiger eines tief eingebetteten Elements im Satz sein. Die Position
vor dem V in den V2-Sprachen hat also nichts mit der SVO/SOV-Dichotomie zu tun; sie stört im Grunde
das Bild und macht den im WALS verfolgten Zählansatz (siehe Abschnitt~\ref{sec-intro-svo})
unanwendbar.\footnote{
  Martin Haspelmath (p.\,M.\ 2022) hat darauf hingewiesen, dass der Zählansatz vergleichenden
  Zwecken dient und nicht deskriptiven. \citet{Haspelmath2010a} -- der zwischen vergleichenden Konzepten und deskriptiven
  Kategorien unterscheidet -- schreibt, dass vergleichende Konzepte nicht richtig oder falsch sein können (S.\,665), sie können nur mehr oder
  weniger für bestimmte Zwecke geeignet sein. Ich stimme zu, dass der Zählansatz nützlich ist, um die
  Sprachen der Welt zu vergleichen, aber mein Punkt war, dass er für die \ili{germanisch}en Sprachen nicht geeignet ist, da hier
  das V2-Phänomen das Bild stört. Zur weiteren Diskussion vergleichender Konzepte vs.\ deskriptiver Kategorien
  siehe \citew{Newmeyer2010a}, \citew{Haspelmath2010b} und \citew[\page 43--44]{MuellerCoreGram}.
}
% \itdopt{M: The counting approach serves comparative purposes, while the Vorfeld
%   fronting approach serves descriptive purposes -- both are justified. Comparative concepts are not
%   appropriate as descriptive devices \citep{Haspelmath2010a}.}

Im Folgenden führe ich Fakten an, die als Evidenz für SOV als Grundabfolge des \ili{Deutsch}en angesehen werden
(und anderer \ili{germanisch}er Sprachen, \zb \ili{Niederländisch}, \ili{Friesisch}, \ili{Afrikaans} und ihrer regionalen
Varianten). Bevor ich in Abschnitt~\ref{sec-analysis-verb-mevement} eine Analyse vorlege, diskutiere ich die
Verbstellung in den \ili{germanisch}en SVO-Sprachen am Beispiel des \ili{Dänisch}en in Abschnitt~\ref{sec-danish-verb-movement}.

\subsection{Deutsch als SOV-Sprache}
\label{sec-German-as-SOV}

\subsubsection{Die Abfolge von Partikel und Verb}

Verbpartikeln\is{Verb!Partikel} bilden eine enge Einheit mit dem Verb. Die Einheit ist nur in verb"=finalen Sätzen beobachtbar,
was eine SOV-Analyse stützt \citep[\page 35]{Bierwisch63a}.
\eal
\ex 
weil er morgen anfängt
\ex 
Er fängt morgen an.
\zl

%\largerpage
\noindent
Das Partikelverb in (\mex{0}) ist nicht transparent: Seine Bedeutung steht in keinem Zusammenhang mit dem Verb
\emph{fangen}. Solche Partikelverben werden manchmal Mini-Idiome
genannt.

\subsubsection{Idiome}

Das obige Argument lässt sich auch mit Idiomen führen, die keine Partikelverben enthalten: Viele Idiome erlauben keine
Umstellung der Idiom\is{Idiom}teile im \mf:\footnote{
  Wie das Beispiel (\ref{ex-niemand-macht}) zeigt, kann das Verb in Anfangsstellung verwendet werden. \citet[\page
  203]{MuellerLehrbuch3} diskutiert Beispiele, in denen die Phrase \emph{den Garaus} im \vf steht.
}
\eal
\judgewidth{?*}
\ex[]{
dass niemand dem Mann den Garaus macht
}
\ex[?*]{
dass dem Mann den Garaus niemand macht
}
\ex[]{
\label{ex-niemand-macht}
Niemand macht ihm den Garaus.
}
\zl
Dies ist ein Beispiel für Behaghels Gesetz \citeyearpar{Behaghel32-u},
dass Dinge, die semantisch zusammengehören, tendenziell zusammen realisiert werden. Die Ausnahme ist das
finite Verb. Das finite Verb kann in Anfangs- oder Endstellung realisiert werden, obwohl dies
die Kontinuität des idiomatischen Materials unterbricht. Da die Kontinuität nur in SOV-Abfolge
beobachtbar ist, wird diese Abfolge als basal angesehen.

\subsubsection{Durch Rückbildung gebildete Verben}

%\largerpage[2]
Verben, die durch Rückbildung\is{Rückbildung|(} aus Substantiven abgeleitet sind, können oft nicht getrennt werden, und verb"=zweite
Sätze sind daher ausgeschlossen (siehe \citealt[\page 62]{Haider93a}, der auf unveröffentlichte Arbeiten
von \citealt{Hoehle91b} verweist, die nun in einer Sammlung von Höhles Arbeiten bei Language Science Press veröffentlicht sind. Die
Beispiele finden sich auf Seite 370--371):
\eal
\label{ex-uraufführen}
\ex[]{
weil sie das Stück heute uraufführen
}
\ex[*]{
Sie uraufführen heute das Stück.
}
\ex[*]{
Sie führen heute das Stück urauf.
}
\zl
Daher können diese Verben nur in der Abfolge verwendet werden, die als Grundabfolge angenommen wird.%
\is{Rückbildung|)}

\subsubsection{Doppelpartikelverben}

Beispiele mit Rückbildung wie in (\ref{ex-uraufführen}) wurden als speziell kritisiert, da sie eine
Rückbildung beinhalten. Vielleicht ist also ein bestimmter, schlecht verstandener Aspekt für ihre Eigenschaften verantwortlich. Aber
\citet[\page 63]{Haider93a},
\citet[\page 105]{Vikner2001a}, \citet{Fortmann2007a} und \citet[\page 59--60]{Haider2010a} haben eine ähnliche
Klasse von Verben gefunden, die sich der Bewegung widersetzen: Doppelpartikelverben. Verben wie \emph{vorankündigen}
bestehen aus der Kombination von \emph{an} und \emph{kündigen} mit der Hinzufügung einer weiteren Partikel \emph{vor}.
\ea 
\label{ex-dass-sie-es-vorankündigt}
dass sie es vor-an-kündigt
\z
Das Interessante an diesen Verben ist nun, dass sie nicht vorangestellt werden können. Der Verbstamm muss
benachbart zur Partikel stehen:
\eal
\ex[*]{
Sie kündigt$_i$ es vor-an \_$_i$.
}
\ex[*]{
Sie an-kündigt$_i$ es vor \_$_i$.
}
\ex[*]{
Sie vor-an-kündigt$_i$ es  \_$_i$.
}
\zl
Die Beispiele zeigen, dass Doppelpartikelverben in OV-Abfolge möglich, in VO-Abfolge aber unmöglich sind.

\subsubsection{Konstruktionen, die nur SOV-Abfolge erlauben}

%\largerpage
Ebenso ist es unmöglich, das Verb in Anfangsstellung zu realisieren, wenn Elemente wie
\emph{mehr als} im Satz vorhanden sind \parencites[Section~3.1]{Haider97c}[\page 732]{Meinunger2001a}:
\eal
\ex[]{
dass Hans seinen Profit letztes Jahr mehr als verdreifachte
}
\ex[]{
Hans hat seinen Profit letztes Jahr mehr als verdreifacht.
}
\ex[*]{
Hans verdreifachte seinen Profit letztes Jahr mehr als.
}
\zl
%\itdopt{AZ: what about Hans mehr als verdreifachte seinen Profit letztes Jahr.}
% surprisingly good but bad.
%
Es ist also möglich, das Adjunkt zusammen mit dem Verb in Endstellung zu realisieren, aber es gibt
Beschränkungen hinsichtlich der Platzierung des finiten Verbs in Anfangsstellung.

\subsubsection{Abfolge in Neben- und infiniten Sätzen}

Verben in infiniten Sätzen und in eingeleiteten finiten Nebensätzen
  erscheinen immer final, das heißt, in der rechten Satzklammer. Zum Beispiel erscheinen \emph{zu geben}
  und \emph{gibt} in (\mex{1}) in der rechten Satzklammer:
%(\mex{1}a) and (\mex{1}b):
\eal
\ex 
Der Clown versucht, Kurt-Martin die Ware zu geben.
\ex 
dass der Clown Kurt-Martin die Ware gibt
\zl

\subsubsection{Skopus von Adverbialen}

%%\largerpage[3]
Der\is{Skopus|(} Skopus von Adverbialen in Sätzen wie (\ref{bsp-absichtlich-nicht-anal}) hängt von ihrer Abfolge ab \citep[Abschnitt~2.3]{Netter92}:
Das am weitesten links stehende Adverb skopiert über das folgende Adverb und über das Verb in
Endstellung. Dies wurde durch die Annahme der folgenden Struktur erklärt:
\eal
\label{bsp-absichtlich-nicht-anal}
\ex 
weil er  [absichtlich [nicht lacht]]
\ex 
weil er [nicht [absichtlich lacht]]
\zl
%\largerpage
Eine interessante Tatsache ist, dass sich die Skopusverhältnisse nicht ändern, wenn die Verbstellung geändert wird. Wenn
man annimmt, dass die Sätze eine zugrundeliegende Struktur wie in (\mex{0}) haben und dass der Skopus
mit Bezug auf diese Struktur bestimmt wird, ist diese Tatsache automatisch erklärt:
\eal
\label{bsp-absichtlich-nicht-anal-v1}
\ex 
Lacht$_i$ er [absichtlich [nicht \_$_i$]]?
\ex 
Lacht$_i$ er [nicht [absichtlich \_$_i$]]?
\zl
%\item Verum-Fokus
\nocite{Hoehle88a,Hoehle97a}

\noindent
Hier ist zu erwähnen, dass es Ausnahmen von der Behauptung zu geben scheint, dass Modifikatoren von
links nach rechts skopieren. \citet*[\page47]{Kasper94a} diskutiert die Beispiele in (\mex{1}), die auf
\citet*[\page137]{BV72} zurückgehen.

\eal
\label{bsp-peter-liest-gut-wegen}
\ex 
Peter liest wegen der Nachhilfestunden gut.
\ex 
Peter liest gut wegen der Nachhilfestunden.
\zl
(\mex{0}a) entspricht der erwarteten Abfolge, in der die adverbiale PP \emph{wegen der
  Nachhilfestunden} über das Adverb \emph{gut} skopiert, aber die alternative Abfolge in (\mex{0}b) ist
ebenfalls möglich, und der Satz hat dieselbe Lesart wie der in (\mex{0}a).

% Kiss95b:212
  \citet[Abschnitt~6]{Koster75a} und \citet*[\page67]{Reis80a} haben jedoch gezeigt, dass diese Beispiele
  keine überzeugende Evidenz sind, da die rechte Satzklammer nicht gefüllt ist und die
  Abfolgen in (\mex{0}) daher nicht notwendigerweise Varianten von \emph{Mittelfeld}abfolgen sind, sondern auf der Extraposition
  einer Konstituente beruhen können. Wie Koster und Reis gezeigt haben, werden die Beispiele ungrammatisch, wenn die rechte
  Satzklammer gefüllt ist:
\eal
\ex[*]{
Hans hat gut wegen der Nachhilfestunden gelesen.
}
\ex[]{
Hans hat gut gelesen wegen der Nachhilfestunden.
}
\zl
Die Schlussfolgerung ist, dass (\mex{-1}b) am besten als Variante von (\mex{-1}a) behandelt wird, in der die PP
extraponiert ist und der Skopus an der Position bestimmt wird, an der die PP stehen würde, wenn sie nicht extraponiert wäre.

Während Beispiele wie (\mex{-1}) zeigen, dass die Sache nicht trivial ist, zeigt das folgende Beispiel von \citet[\page
383]{Crysmann2004a}, dass es Beispiele mit gefüllter rechter Satzklammer gibt, die
Skopusverhältnisse erlauben, in denen ein Adjunkt über ein anderes, vorangehendes Adjunkt skopiert. So skopiert zum Beispiel in
(\mex{1}) \emph{niemals} über \emph{wegen schlechten Wetters}:
\ea
 Da muß es schon erhebliche Probleme mit der Ausrüstung gegeben haben, da [wegen
  schlechten  Wetters] ein Reinhold Messner [niemals] aufgäbe.
%\ex Stefan  ist wohl deshalb krank geworden, weil er äußerst hart wegen der Konferenz in Bremen gearbeitet hat.
\z

\noindent
Dies ändert jedoch nichts an der Tatsache, dass die Sätze in (\ref{bsp-absichtlich-nicht-anal}) und
(\ref{bsp-absichtlich-nicht-anal-v1}) unabhängig von der Stellung des Verbs dieselbe Bedeutung haben. Die allgemeine
Bedeutungskomposition kann auf die Weise erfolgen, die Crysmann vorgeschlagen hat.%
%\later{T: i-scrambling}
%

Ein weiteres Wort der Vorsicht ist hier angebracht: Es gibt SVO-Sprachen wie das \ili{Französisch}e, die ebenfalls eine
Links-nach-rechts-Skopierung von Adjunkten aufweisen \citep[\page 156--161]{BGK2004a-u}. Die obige Argumentation sollte daher nicht als der einzige
Fakt angesehen werden, der den SOV-Status des \ili{Deutsch}en stützt. In jedem Fall können die in verschiedenen Rahmentheorien
ausgearbeiteten Analysen des \ili{Deutsch}en die Fakten gut erklären.
\is{Skopus|)}

\subsubsection{Stellung infiniter Verben in VO- und OV-Sprachen}

Bevor ich mich im nächsten Unterabschnitt der Verbstellung im \ili{Dänisch}en zuwende, möchte ich Ørsnes'
Beispiele mit mehreren infiniten Verben wiederholen (siehe (\ref{ex-Danish-embedding}) auf S.\,\pageref{ex-Danish-embedding}): Das Beispiel in (\mex{1}a) zeigt einen \ili{deutsch}en Nebensatz mit einem Verbalkomplex aus
drei Verben. Die Einbettungsebene ist durch tiefgestellte Zahlen angegeben. Wie man sieht, werden die Verben
am Ende des Satzes angefügt. Im entsprechenden \ili{dänisch}en Beispiel, das aus \citet[\page 146]{Oersnes2009b} adaptiert wurde, ist es genau umgekehrt:
Das einbettende Verb geht dem eingebetteten Verb voran.
\eal
\ex
 dass er ihn gesehen$_3$ haben$_2$ muss$_1$
\ex
\gll at hun må$_1$ have$_2$ set$_3$ ham\\
     dass sie muss haben gesehen ihn\\\danish
\glt `dass sie ihn gesehen haben muss'
\zl
%
Die Beispiele in (\mex{1}) sind Varianten mit unterschiedlicher Komplexität. Wenn wir das Verb
\emph{sah} in (\mex{1}a) durch die Perfektform ersetzen, wird das Auxiliar nach dem Partizip platziert wie
in (\mex{1}b).
\eal
\ex
dass er ihn sah
\ex
dass er ihn gesehen hat
\zl 
Wird zu (\mex{0}b) ein Modalverb hinzugefügt, so steht das Modalverb rechts von den eingebetteten Verben. Diese Abfolge wird
durch die Platzierung des finiten Verbs in Anfangsstellung verzerrt, aber diese Platzierung ist unabhängig
von der Abfolge der infiniten Verben. Wie die Beispiele in (\mex{1}) zeigen, wird das finite Verb
sowohl im \ili{Deutsch}en (SOV) als auch im \ili{Dänisch}en (SVO) links vom Subjekt realisiert.

\eal
\ex 
 Muss er ihn gesehen haben?
\ex 
\gll Må han have set ham?\\
     muss er ihn gesehen haben\\\danish
\glt `Muss er ihn gesehen haben?'
\zl
\is{Abfolge!SOV|)}

\subsection{Verbstellung in den germanischen SVO-Sprachen}
\label{sec-danish-verb-movement}\label{sec-Germanic-SVO-verb-position}

Während der Diskussion der Skopusfakten habe ich bereits eine Analyse angedeutet, in der eine \isi{Spur} die
Position des Verbs in Endstellung markiert und das Verb in Anfangsstellung mit dieser
Spur koindiziert ist. Obwohl die SVO-Sprachen verschieden sind, wurde eine ähnliche Analyse für Sprachen
wie das \ili{Dänisch}e vorgeschlagen. Die Evidenz dafür ist, dass sich Adverbiale in SVO-Sprachen üblicherweise an die VP anschließen,
das heißt, sie verbinden sich mit einer Phrase, die aus dem Verb und seinem Objekt oder seinen Objekten besteht. (\mex{1}) ist
ein Beispiel:
\ea
\gll  at   Conny ikke [\sub{VP} læser bogen]\\
      dass Conny nicht      {}        liest          Buch.\textsc{def}\\\danish
\glt `dass Conny das Buch nicht liest'
\z

Das Interessante daran ist nun, dass das finite Verb in V2-Sätzen links von der Negation platziert wird:

\ea
\gll  Conny læser ikke bogen.\\
       Conny liest   nicht  Buch.\textsc{def}\\\danish
\glt `Conny liest das Buch nicht.'
\z
Dies wird von vielen als Evidenz für Verbvoranstellung angesehen:
\ea
\gll  Conny læser$_i$ ikke [\sub{VP} \_$_i$ bogen].\\
      Conny liest      nicht  {} {}    Buch.\textsc{def}\\\danish
\glt `Conny liest das Buch nicht.'
\z
\nocite{KS2002a}

Vor diesem Hintergrund sollte klar sein, wie die Analyse von Entscheidungsfragen wie in (\mex{1}b) aussieht:
\eal
\ex
\gll at Conny læser bogen\\
     dass Conny liest Buch.\textsc{def}\\\danish
\glt `dass Conny das Buch liest'
\ex\label{ex-laeser-jens-bogen}
\gll Læser Conny bogen?\\
     liest Conny Buch.\textsc{def}\\
\glt `Liest Conny das Buch?'
\zl
Die Analyse des ersten Satzes enthält eine VP wie in (\mex{1}a), und der zweite Satz enthält eine
VP mit einer verbalen \isi{Spur}, die dem Verb in Anfangsstellung entspricht:
\eal
\ex
\gll at Conny [\sub{VP} læser bogen]\\
     dass Conny {} liest Buch.\textsc{def}\\\danish
\glt `dass Conny das Buch liest'

\ex
\gll Læser$_i$ Conny [\sub{VP} \_$_i$ bogen]?\\
     liest     Conny {}        {}     Buch.\textsc{def}\\
\glt `Liest Conny das Buch?'
\zl

Es ist interessant festzustellen, dass die \ili{deutsch}e und die \ili{dänisch}e Frage mit nur einem einzigen Verb genau
dieselbe Konstituentenabfolge haben. Man vergleiche (\ref{ex-laeser-jens-bogen}) mit (\mex{1}):
\ea
Liest Conny das Buch?
\z
Die interne Struktur dieser Sätze ist jedoch recht unterschiedlich. Der unterschiedliche Charakter der beiden
Sprachen wird natürlich offensichtlicher, wenn infinite Verben beteiligt sind:
\eal
\ex
\gll Har$_i$ Conny [ \_$_i$ læst bogen]?\\
     hat Conny {} {} gelesen Buch.\textsc{def}\\\danish
\glt `Hat Conny das Buch gelesen?'
\ex
Hat$_i$ Conny das Buch [gelesen \_$_i$]?
\zl
In (\mex{0}a) ist das finite Verb mit einer \isi{Spur} in Anfangsstellung der VP verbunden und in
(\mex{0}b) ist es mit einem Verb in Endstellung in einem Verbalkomplex verbunden.%
\is{Abfolge!SVO|)}

\subsection{Verb second}
\label{sce-verb-second}

Sogar\is{Abfolge!V2|(} Sprachen mit recht starrer Konstituentenabfolge erlauben manchmal das Voranstellen von Elementen.
% AZ: leave out or provide example
% or positioning of elements at the far right, that is, their extraposition.
(\mex{1}) zeigt \ili{englisch}e Beispiele für Voranstellung:
\eal
\ex I read this book yesterday.
\ex This book, I read yesterday.
\ex Yesterday, I read this book.
\zl
Das Objekt \emph{this book} und das Adjunkt \emph{yesterday} sind in (\mex{0}b) bzw.\
(\mex{0}c) vorangestellt.

Die \ili{germanisch}en Sprachen (mit Ausnahme des \ili{Englisch}en) stellen eine Konstituente vor das finite
Verb. Wie die \ili{deutsch}en Beispiele in (\mex{1}) zeigen, kann die vorangestellte Konstituente eine beliebige grammatische Funktion haben:
\eal
\ex 
Ich habe das Buch gestern gelesen.
\ex 
Das Buch habe ich gestern gelesen.
\ex 
Gestern habe ich das Buch gelesen.
\ex 
Gelesen habe ich das Buch gestern, gekauft hatte ich es aber schon vor einem Monat.
\ex 
Das Buch gelesen habe ich gestern.
\zl
Solche Voranstellungen sind nicht satzgebunden, das heißt, die Voranstellung kann eine oder mehrere Satzgrenzen
und auch Grenzen anderer Konstituenten überschreiten. (\mex{1}) zeigt \ili{englisch}e Beispiele, in denen das Objekt von
\emph{saw} über eine und über zwei Satzgrenzen hinweg extrahiert ist:
\eal
\ex\label{ex-chris-we-saw} Chris, Sandy saw.
\ex\label{ex-chris-we-think-that-sandy-saw} Chris, we think that Sandy saw.
\ex Chris, we think Anna claims that Sandy saw.
\zl
Im \ili{Deutsch}en finden sich solche Extraktionen ebenfalls:
\eal
\label{ex-fernabhaengigkeit-one}
\ex 
 Wer$_i$ wohl    meint   er, dass \_$_i$ ihm seine Arbeit hier bezahlen werde?\footnote{
\citew[\page 321]{Paul1919a}. Paul provides two pages full of attested examples of extractions out of
\emph{dass} clauses.
}

\ex
\label{ex-wen-glaubst-du-dass}
 Wen$_i$ glaubst du, daß ich \_$_i$ gesehen habe.\footnote{
    \citew[\page84]{Scherpenisse86a}.
    }

\ex 
 "`Wer$_i$, glaubt\iw{glauben} er, daß er \_$_i$ ist?"' erregte sich ein Politiker vom Nil.\footnote{
        Spiegel, 1999-08, p.\,18.
}

\zl
Man sagt im Allgemeinen, dass sie in süddeutschen Varietäten häufiger sind, aber es gibt andere
Beispiele, die zeigen, dass Fernabhängigkeiten beteiligt sind. In (\ref{bsp-um-zwei-millionen}) hängt das
präpositionale Objekt \emph{um zwei Millionen Mark} von
\emph{betrügen} ab. Es hängt
von keinem der Verben im Matrixsatz ab. Die Phrase \emph{eine Versicherung zu betrügen} ist
extraponiert, das heißt, sie steht rechts von der Verbalklammer im sogenannten \nf. Die
Position von \emph{um zwei Millionen Mark} kann nicht durch lokale Umstellung erklärt werden. Ebenso hängt
\emph{gegen ihn} von \emph{Angriffe} ab, das Teil der Phrase \emph{Angriffe zu
  lancieren} ist. Wieder ist eine Analyse, die auf lokaler Umstellung von Abhängigen eines Kopfes beruht, unmöglich.
\eal
\label{bsp-Fernabhaengigkeit}
\ex\label{bsp-um-zwei-millionen}
 {}[Um zwei Millionen Mark]$_i$ soll er versucht haben, [eine Versicherung \_$_i$ zu betrügen].\footnote{
         taz, 2001-05-04, p.\,20.
}

\ex
 {}[Gegen ihn]$_i$ falle es den Republikanern hingegen schwerer, [~[~Angriffe \_$_i$] zu lancieren].\footnote{
  taz, 2008-02-08, p.\,9.
}

\zl
\is{Abfolge!V2|)}

\section{Die Analyse}
\label{sec-analysis-verb-mevement}

Dieser Abschnitt behandelt verb"=erste Sätze in Unterabschnitt~\ref{sec-verb-first} und danach
verbzweite Sätze in Unterabschnitt~\ref{sec-verb-second}. Ich erläutere zwei verschiedene Mechanismen für die
jeweiligen Analysen: die Doppelschrägstrich-Notation für Abhängigkeiten, die einen Kopf betreffen (sogenannte „Kopfbewegung“), und die Schrägstrich-Notation
für Fernabhängigkeiten (sogenannte „Konstituentenbewegung“).\footnote{
  Man beachte, dass sich HPSG von Theorien wie Government \& Binding (GB, \citealt{Chomsky81a}) und
  Minimalismus \citep{Chomsky95a-u} darin unterscheidet, dass nichts
  tatsächlich bewegt wird. Während GB annimmt, dass es zwei oder mehr Strukturen gibt, die durch
  Bewegung von Konstituenten miteinander in Beziehung stehen, nimmt HPSG nur eine einzige Struktur mit einem leeren Element und Informationsteilung
  an. Dies ist ein wichtiger Unterschied, was die psycholinguistische Plausibilität von Theorien
  betrifft. Siehe \citew[Chapter~15]{MuellerGT-Eng5} und \citew[Section~3.2]{Wasow2021a} zur Diskussion.
}
Die Analyse der „Kopfbewegung“ geht auf \citet{Jacobson87} zurück, der eine solche Analyse
im Rahmen der Kategorialgrammatik für das \ili{Englisch}e vorgeschlagen hat. \citet{Borsley89} hat diese Analyse für HPSG adaptiert,
und \citet{KW91a,Kiss95a,Kiss95b} haben eine Verbbewegungsanalyse für das \ili{Deutsch}e in HPSG vorgeschlagen. Die Analyse der
Konstituentenbewegung ist tatsächlich älter als die der Kopfbewegung. Sie geht auf Arbeiten in der
Generalized Phrase Structure Grammar \citep{Gazdar81a} zurück und wurde von \citet[Abschnitt~3.4]{ps} und \citet[Kapitel~4]{ps2} für HPSG adaptiert.

\subsection{Verberststellung}
\label{sec-verb-first}

%\largerpage
Die\is{Kopfbewegung|(} Analyse verwendet einen speziellen Mechanismus, der Informationen in einem Baum nach oben reicht. Die
Verbspur\is{Spur!Kopfbewegung|(} enthält die
Information, dass lokal ein Verb fehlt. Diese Information über das fehlende Verb wird an den
Knoten weitergegeben, der die Verbspur dominiert. Sie wird mithilfe eines Mechanismus dargestellt, der „Doppelschrägstrich“
genannt und als //\is{\dslmath} geschrieben wird. Die
betreffende Information ist Kopfinformation und wird daher entlang des Kopfpfades zusammen mit
anderen Informationen, etwa der Wortart, nach oben gereicht. Abbildung~\vref{fig-analysis-German-verb-first}
illustriert das. Der Verbspur fehlt ein V, dem V$'$ fehlt ein V, und dem S ebenfalls.
\begin{figure}
\centering
\begin{forest}
sm edges
[S
  [{V \sliste{ S//V }}
    [V [liest$_j$] ] ]
       [{S//V}
           [NP [Conny] ]
           [\vbarDSLv
             [NP [das Buch, roof] ]
             [{\mybox[v1]{V}//\mybox[v2]{V}} [\_$_j$] ] ] ] ] ]
%\draw[semithick,<->,color=green] (v1.south)--(v2.south);
%% \draw[semithick,<->,color=green] (3.1,-3.9) ..controls +(south east:.5) and +(south west:.5)..(2.7,-3.9);
%% \draw[semithick,<->,color=green] (3.5,-3.7) ..controls +(east:.5) and +(east:.5)..(2.8,-2.5);
%% \draw[semithick,<->,color=green] (2.8,-2.3) ..controls +(east:.5) and +(east:.5)..(1.7,-1.1);
%% \draw[semithick,<->,color=green] (1.5,-0.9) ..controls +(north:.5) and +(north:.5)..(-0.8,-0.9);
%% \draw[semithick,<->,color=green] (-0.7,-1.1) ..controls +(south east:.2) and +(north
       %% east:.5)..(-1.0,-2.4);
\end{forest}
\caption{\label{fig-analysis-German-verb-first}\label{fig-liest-jens-das-buch}Analyse der Verbstellung im Deutschen}
\end{figure}
Das initiale Verb selegiert einen Satz, dem ein V fehlt \sliste{ S//V }. Der Lexikoneintrag für
das Verb in Anfangsstellung wird durch eine Lexikonregel lizenziert, die ein Verb mit einem Verb in Beziehung setzt, das
einen Satz selegiert, dem das Eingabeverb fehlt. Da die Selektionsforderung dieses
Verbs (S//V) mit dem Satz, dem ein V fehlt (Conny das Buch \_$_j$), identifiziert wird, wird die Information
über das ursprüngliche Verb \emph{liest} mit dem V in S//V identifiziert. Da die Doppelschrägstrich-Information
Kopfinformation ist, perkoliert sie entlang des Kopfpfades nach unten zur Verbspur. Die
Information über das initiale V wird mit der syntaktischen und semantischen Information der
Verbspur in Endstellung identifiziert, und daher verhält sich diese Verbspur genau wie das Verb in Anfangsstellung,
das die Eingabe der Lexikonregel war.

Verschiedene Forscherinnen und Forscher haben argumentiert, dass sich das finite Verb in Anfangsstellung wie ein Komplementierer in
Nebensätzen verhält \citep{Hoehle97a,Weiss2005a-u,Weiss2018a-u}. Dies wird von der Analyse erfasst. Man vergleiche
Abbildung~\ref{fig-analysis-German-verb-first} mit Abbildung~\vref{fig-analysis-German-verb-last}.
%\largerpage
\begin{figure}
\centering
\begin{forest}
sm edges
[CP
  [{C \sliste{ S }} [dass] ]
  [S
           [NP [Conny] ]
           [V$'$
             [NP [das Buch, roof] ]
             [V [liest] ] ] ] ]
\end{forest}
\caption{\label{fig-analysis-German-verb-last}Analyse eines verb"=finalen Satzes mit Komplementierer im Deutschen}
\end{figure}
Der Komplementierer \emph{dass} selegiert einen vollständigen Satz, das heißt, einen Satz, dem
kein Verb fehlt, und das initiale Verb \emph{liest} in Abbildung~\ref{fig-analysis-German-verb-first} selegiert
einen Satz, dem \emph{liest} fehlt. Abgesehen vom overten oder kovertem Verb sind die Strukturen
also identisch. Diese Tatsache ist wichtig, wenn es um die Analyse der Skopusfakten geht.
\begin{figure}
\centering
\begin{forest}
sm edges
[S
        [{V \sliste{ S//V }} 
          [V [lacht$_j$] ] ]
        [{S//V}
           [NP [er] ]
           [\vbarDSLv
             [Adv [nicht] ]
             [\vbarDSLv
               [Adv [absichtlich] ]
               [{V//V} [\_$_j$] ] ] ] ] ]
\end{forest}
\caption{\label{fig-analysis-German-verb-initial-scope}Analyse von Sätzen mit Adverbialen im Deutschen}
\end{figure}
Da die Struktur völlig parallel zu der ist, die wir in verb"=finalen Sätzen haben, folgen die Skopusfakten
unmittelbar: Die Spur verhält sich wie das Verb in Anfangsstellung, \emph{absichtlich}
modifiziert die Spur, und die resultierende Semantik wird im Baum nach oben gereicht (siehe Abbildung~\vref{fig-analysis-German-verb-initial-scope}). Der nächste
Schritt ist die Modifikation durch \emph{nicht}. Wieder wird die resultierende Semantik nach oben
gereicht. \emph{lacht} verbindet sich mit dem Satz und übernimmt dessen Semantik. Da \emph{lacht}
der Kopf ist, wird die Semantik von dort aus weitergegeben.

Die Analyse des \ili{Dänisch}en ist völlig parallel zur Analyse des \ili{Deutsch}en. Der einzige Unterschied zwischen
Abbildung~\ref{fig-analysis-German-verb-first} und Abbildung~\vref{fig-analysis-verb-first-Danish} ist die
Position der Verbspur relativ zum Objekt: Die Spur folgt dem Objekt im \ili{Deutsch}en, aber
sie geht ihm im \ili{Dänisch}en voran.
\begin{figure}
\centering
\begin{forest}
sm edges
[S
        [{V \sliste{ S//V }} 
          [V [læser$_j$;liest] ] ]
        [{S//V}
           [NP [Conny;Conny] ]
           [{VP//V}
             [{V//V} [\_$_j$] ]
             [NP [bogen;Buch.\textsc{def}] ] ] ] ]
\end{forest}
\caption{\label{fig-analysis-verb-first-Danish}Analyse der Verbstellung im Dänischen}
\end{figure}%

Das Letzte, was in diesem Abschnitt erklärt wird, ist die Analyse der Negation und der Verbvoranstellung im
\ili{Dänisch}en. Abbildung~\vref{fig-analysis-verb-fronting-negation-Danish} zeigt, dass sich die Negation an
die VP anschließt wie in verb"=finalen Sätzen und das Verb vorangestellt wird, sodass es links von der Negation erscheint.
\begin{figure}
\centering
\begin{forest}
sm edges
[S
        [{V \sliste{ S//V }} 
          [V [læser$_j$;liest] ] ]
        [{S//V}
           [NP [Conny;Conny] ]
           [{VP//V}
             [Adv [ikke;nicht] ]
             [{VP//V}
               [{V//V} [\_$_j$] ]
               [NP [bogen;Buch.\textsc{def}] ] ] ] ] ]
\end{forest}
\caption{\label{fig-analysis-verb-fronting-negation-Danish}Die Analyse von Verbvoranstellung und
  Negation im Dänischen}
\end{figure}
Der nächste Abschnitt erklärt die Extraktion von Konstituenten, und es wird dann möglich sein, die
vollständige Struktur für Sätze wie (\mex{1}a) anzugeben. Es wird auch klar werden, warum die Abfolge
von Negation und Verb sich in eingebetteten und in Hauptsätzen unterscheidet:
\eal
\ex
\gll Conny læser ikke bogen.\\
     Conny liest nicht  Buch.\defsc\\
\glt `Conny liest das Buch nicht.'
\ex
\gll at Conny ikke læser bogen\\
     dass Conny nicht liest Buch.\defsc\\
\glt `dass Conny das Buch nicht liest'
\zl
\is{Spur!Kopfbewegung|)}\is{Kopfbewegung|)}

\subsection{Verbzweitstellung}
\label{sec-verb-second}

Die\is{Vorfeldbesetzung|(} Technik, die für die Analyse von Fernabhängigkeiten verwendet wird, ist dieselbe, die
für die Analyse der Umstellungen von Verben eingesetzt wurde: Ein leeres Element nimmt die Position der vorangestellten
Konstituente ein, und die Information über die fehlende Konstituente (die sogenannte \emph{Lücke}) wird im
Baum nach oben gereicht, bis sie schließlich durch das vorangestellte Element abgebunden wird, den sogenannten
\emph{\isi{Füller}}. Abbildung~\vref{fig-chris-sandy-saw} illustriert die Analyse von (\ref{ex-chris-we-saw}).

\begin{figure}
\begin{forest}
sm edges
[S
  [NP [Chris] ]
  [S/NP 
    [NP [Sandy] ] 
    [VP/NP  
      [V [saw] ]
      [NP/NP [\trace] ] ] ] ]
%% \draw[connect] (NP/NP.north east) [bend right] to (VP/NP.south east);
%% \draw[connect] (VP/NP.north east) [bend right] to (S/NP.south east);
%% \draw[connect] (S/NP.north east) [bend right] to (NP);
\end{forest}
\caption{\label{fig-chris-sandy-saw}Die Analyse der Extraktion im Englischen}
\end{figure}
Die Kategorie, die auf den Schrägstrich folgt (`/'\is{/}), steht für das Objekt, das lokal in der Position
der Spur fehlt. Spuren sind wie Joker in Kartenspielen: Sie können (fast) jede Position füllen. Sie
geben vor, von der Kategorie zu sein, die lokal verlangt wird (die NP im Akkusativ im vorliegenden
Beispiel), aber die Information, dass diese Kategorie lokal fehlt, wird nach oben gereicht (von NP/NP zu VP/NP
zu S/NP). Wenn ein passender Füller mit einer Schrägstrich-Konstituente verbunden wird, wird die Information über das fehlende Element nicht
weiter nach oben gereicht. Man sagt, dass die Fernabhängigkeit an diesem Punkt abgebunden wird. In
Abbildung~\ref{fig-chris-sandy-saw} wird S/NP mit dem Füller NP verbunden, und daher ist der Mutterknoten ein
S ohne Schrägstrich. Das Verb hat alle seine Argumente, und kein Schrägstrich-Element fehlt im Satz: Der
Satz ist vollständig.

%\largerpage
Abbildung~\vref{fig-chris-we-think-that-sandy-saw} zeigt die Analyse von Beispiel
(\ref{ex-chris-we-think-that-sandy-saw}), das tatsächlich eine Fernabhängigkeit erfordert. Wie in
der Abbildung gezeigt, wird die Information über das fehlende Objekt auf die Satzebene (S/NP),
auf die CP-Ebene (CP/NP) und bis zum nächsthöheren S nach oben gereicht. Dort wird sie durch den Füller \emph{Chris} abgebunden.
\begin{figure}
\scalebox{.95}{%
\begin{forest}
sm edges
[S
  [NP [Chris] ]
  [S/NP
    [NP [we] ] 
    [VP/NP  
       [V [think] ]
       [CP/NP
         [C [that] ]
         [S/NP
            [NP [Sandy] ] 
            [VP/NP  
               [V [saw] ]
               [NP/NP [\trace ] ] ] ] ] ] ] ]
%% \draw[connect] (NP/NP.north east)  [bend right] to (VP/NP.south east);
%% \draw[connect] (VP/NP.north east)  [bend right] to (S/NP.south east);
%% \draw[connect] (S/NP.north east)   [bend right] to (CP/NP.south east);
%% \draw[connect] (CP/NP.north east)  [bend right] to (VP/NP1.south east);
%% \draw[connect] (VP/NP1.north east) [bend right] to (S/NP1.south east);
%% \draw[connect] (S/NP1.north east)  [bend right] to (NP);
\end{forest}}
\caption{\label{fig-chris-we-think-that-sandy-saw}Extraktion über die Satzgrenze hinweg}
\end{figure}
Das Abbinden des fehlenden Elements wird durch ein spezielles Schema lizenziert, das
Füller"=Kopf-Schema genannt wird. Abbildung~\vref{fig-filler-head} gibt eine Skizze dieses Schemas.
\begin{figure}
\begin{forest}
[{S[\type{fin}]}
  [\ibox{1}]
  [{S[\type{fin}]}/\ibox{1}]]
\end{forest}
\caption{\label{fig-filler-head}Skizze des Füller-Kopf\is{Schema!Füller-Kopf}-Schemas}
\end{figure}

Das \ili{Englisch}e ist die einzige Nicht-V2-Sprache unter den \ili{germanisch}en Sprachen. Im Folgenden zeige ich, wie das \ili{Deutsch}e
(V2+SOV) und das \ili{Dänisch}e (V2+SVO) mit den bisher eingeführten Techniken analysiert werden können.
Abbildung~\vref{fig-das-buch-liest-jens} zeigt die Analyse von (\mex{1}):
%\largerpage
\ea
Das Buch liest Conny.
\z
\begin{figure}
\begin{forest}
sm edges
[S
  [NP$_i$ [das Buch, roof] ]
  [S/NP
     [V \sliste{ S//V }
        [V [liest$_j$] ] ]
     [S//\vSLASHnp
        [NP/NP [\trace$_i$] ]
        [\vbarDSLv
           [NP [Conny] ]
           [V//V [\_$_j$] ] ] ] ] ] ]
%% \draw[connect] (NP/NP) [bend right] to (S//\vSLASHnp.south east);
%% \draw[connect] (S//\vSLASHnp.north east) [bend right] to (S/NP.east);
%% \draw[connect] (S/NP.north east) [bend right] to (NP);
\end{forest}
\caption{\label{fig-das-buch-liest-jens}Analyse von V2 im Deutschen (SOV)}
\end{figure}
Die Analyse des \ili{deutsch}en Beispiels ist komplizierter als die des \ili{englisch}en, da Verbbewegung
beteiligt ist. Das Verb wird vorangestellt, wie mit Bezug auf
Abbildung~\ref{fig-liest-jens-das-buch} erläutert wurde. Zusätzlich wird das Objekt durch eine Spur realisiert und dann durch
den Füller \emph{das Buch} gefüllt, der präverbal realisiert wird.

Ich folge \citet{Fanselow2003d} und \citet{Frey2004a}, die annehmen, dass die Position des Objekts
initial im \mf ist. Da das \ili{Deutsch}e sowohl Nominativ-Akkusativ- als auch Akkusativ-Nominativ-Abfolge
erlaubt, könnte die Position der Spur für das extrahierte Objekt initial oder final sein, wie in (\mex{1}a)
bzw.\ (\mex{1}b):
\eal
\ex 
{}[Das Buch]$_i$ liest$_j$ \_$_i$ Conny \_$_j$.
\ex 
{}[Das Buch]$_i$ liest$_j$ Conny \_$_i$ \_$_j$.
\zl
Fanselow und Frey argumentieren, dass vorangestellte Elemente wie \emph{das Buch} informationsstrukturelle Eigenschaften haben,
die denen von nicht-vorangestellten Elementen in der initialen \mf-Position entsprechen:
\ea
Liest das Buch Conny?
\z
Sie argumentieren, dass (\mex{0}) sich wie (\mex{-1}a) und nicht wie (\mex{-1}b) verhält.
Die vollständige Diskussion wird hier nicht wiederholt, da uns dies zu weit wegführen würde, aber die
interessierte Leserin oder der interessierte Leser kann die Diskussion in Abschnitt~\ref{sec-discussion-scope} zurate ziehen.

%\largerpage
Die Analyse des parallelen \ili{dänisch}en V2-Beispiels in (\mex{1}) ist ähnlich.
\ea
\gll Bogen             læser Conny.\\
     Buch.\textsc{def} liest Conny\\
\glt `Conny liest das Buch.'
\z
Die Analyse besteht aus zwei Teilen: erstens der Analyse der Verb-Anfangsstellung, die den
Doppelschrägstrich-Mechanismus involviert, und zweitens der Voranstellung des Objekts mithilfe des Schrägstrich-Mechanismus.
Abbildung~\vref{fig-bogen-laeser-jens} illustriert das.
\begin{figure}
\scalebox{.95}{%
\begin{forest}
sm edges
[S
   [NP$_i$ [bogen;Buch.\textsc{def}] ]
      [S/NP
         [V \sliste{ S//V }
           [V [læser$_j$;liest] ] ]
           [S//\vSLASHnp
             [NP [Conny;Conny] ]
             [VP//\vSLASHnp
               [V//V  [\_$_j$] ]
               [NP/NP [\trace$_i$ ] ] ] ] ] ] ]
%% \draw[connect] (NP/NP.north east) [bend right] to (V/V.south east);
%% \draw[connect] (V/V.north east) [bend right] to (S//\vSLASHnp.south east);
%% \draw[connect] (S//\vSLASHnp.north east) [bend right] to (S/NP.east);
%% \draw[connect] (S/NP.north east) [bend right] to (NP);
\end{forest}}
\caption{\label{fig-bogen-laeser-jens}Analyse von V2 im Dänischen (SVO)}
\end{figure}

Die aufmerksame Leserin oder der aufmerksame Leser wird fragen, warum wir zwei verschiedene Mechanismen verwenden, um Verbbewegung und
Extraktion zu analysieren. Die Antwort ist, dass diese Bewegungstypen von unterschiedlicher Natur sind: Verbbewegung ist
satzgebunden, während die Bewegung anderer Konstituenten Satzgrenzen überschreiten kann. Dies wird durch die
Tatsache erfasst, dass die Doppelschrägstrich-Information zusammen mit anderen Kopfmerkmalen, etwa der
Wortart-Information, nach oben gereicht wird, und die Schrägstrich-Information separat nach oben gereicht wird.

%\largerpage
Bevor wir uns im nächsten Kapitel mit dem Passiv befassen, können wir die drei Sätze in (\mex{1}) vergleichen:
\eal
\ex Conny reads a book.
\ex Conny læser en bog.
\ex Conny liest ein Buch.
\zl
Wieder ist die Abfolge der Elemente in allen drei Sprachen gleich. Allerdings ist das \ili{Englisch}e eine SVO-Sprache
ohne V2, das \ili{Dänisch}e eine SVO+V2-Sprache und das \ili{Deutsch}e eine SOV+V2-Sprache. Die Analysen in
Klammernotation sind in (\mex{1}) angegeben, die Baumstrukturen sind in Abbildung~\vref{fig-English-Danish-German-declerative-main-clause} dargestellt:
\eal
\ex {}[\sub{S} Conny [\sub{VP} reads [\sub{NP} a book]]].
\ex {}[\sub{S} Conny$_i$ [\sub{S/NP} læser$_j$ [\sub{S/NP} \_$_i$ [\sub{VP}  \_$_j$ [\sub{NP} en bog]]]].
\ex {}[\sub{S} Conny$_i$ [\sub{S/NP} liest$_j$ [\sub{S/NP} \_$_i$ [\sub{\vbar} [\sub{NP} ein Buch] \_$_j$]]]].
\zl
\begin{figure}
\scalebox{.8}{%
\begin{forest}
sm edges
[S
  [NP [Conny]]
  [VP
    [V [reads]]
    [NP [a book,baseline,roof]]]]
\end{forest}}
\hfill
\scalebox{.8}{%
\begin{forest}
sm edges
[S
   [NP$_i$ [Conny;Conny] ]
      [S/NP
         [V \sliste{ S//V }
           [V [læser$_j$;liest] ] ]
           [S//\vSLASHnp
             [NP/NP [\trace$_i$ ] ]
             [VP//V
               [V//V  [\_$_j$] ]
               [NP [en bog;ein Buch,baseline,roof] ] ] ] ] ]
%% \draw[connect] (NP/NP.north east) [bend right] to (V/V.south east);
%% \draw[connect] (V/V.north east) [bend right] to (S//\vSLASHnp.south east);
%% \draw[connect] (S//\vSLASHnp.north east) [bend right] to (S/NP.east);
%% \draw[connect] (S/NP.north east) [bend right] to (NP);
\end{forest}}
\hfill
\scalebox{.8}{%
\begin{forest}
sm edges
[S
   [NP$_i$ [Conny] ]
      [S/NP
         [V \sliste{ S//V }
           [V [liest$_j$] ] ]
           [S//\vSLASHnp
             [NP/NP [\trace$_i$ ] ]
             [\vbarDSLv
               [NP [ein Buch,baseline,roof] ]
               [V//V  [\_$_j$] ] ] ] ] ]
\end{forest}}
\caption{\label{fig-English-Danish-German-declerative-main-clause}Deklarative Hauptsätze mit
  Subjekt in Anfangsstellung im \ili{Englisch}en, \ili{Dänisch}en und \ili{Deutsch}en: trotz ähnlicher Erscheinung ist die
  syntaktische Struktur verschieden}
\end{figure}

\noindent
Es mag überraschend sein, dass diese drei Sätze derart radikal verschiedene Analysen erhalten, obwohl die
Abfolge der Elemente gleich ist. Der Unterschied in den Strukturen ist das Ergebnis der Annahme, dass
alle deklarativen Hauptsätze in den \ili{germanisch}en V2-Sprachen demselben Muster folgen, nämlich dass das
finite Verb vorangestellt wird und dann eine weitere Konstituente vorangestellt wird. Diese spezielle Konstruktion ist
mit dem Satztyp verbunden, das heißt, mit der Bedeutung der Äußerung (Imperativ, Frage,
Assertion). Die Sätze in (\ref{ex-fernabhaengigkeit-one}) und (\ref{bsp-Fernabhaengigkeit}) zeigen,
dass V2 eine Fernabhängigkeit involviert. Daher ist die Analyse
von (\mex{0}b) komplexer als (\mex{1}) und involviert die Voranstellung des finiten Verbs in die Anfangsstellung
mit einer anschließenden Voranstellung des Subjekts:
\ea
{}[\sub{S} Conny [\sub{VP} læser [\sub{NP} en bog]]].
\z
%\largerpage
Der Grund ist, dass nun alle deklarativen Hauptsätze unter derselben Struktur subsumiert werden, nämlich
(\mex{-1}b). Ein deklarativer Hauptsatz in allen \ili{germanisch}en V2-Sprachen ist die Kombination einer extrahierten
Phrase mit einer verb"=initialen Phrase, in der das extrahierte Element fehlt. Die Voranstellung des finiten
Verbs ist eine Möglichkeit, den Satztyp zu markieren: Wird nur das finite Verb vorangestellt, ist das Ergebnis eine Entscheidungsfrage (\mex{1}a)
oder ein Imperativsatz (\mex{1}b).\footnote{
  Verb"=initiale Sätze können auch deklarative Sätze sein, wenn sogenannter \emph{Topik-Drop} \citep{Fries88b} beteiligt ist:
  \ea
  Was macht Peter? Gibt ihm ein Buch.
  \z
  Das Subjekt von \emph{gibt} ist weggelassen. Der vollständige Satz wäre ein V2-Satz:
  \emph{Er gibt ihm ein Buch}.
}
\eal
\ex 
Gibt er ihm das Buch?
\ex 
Gib mir das Buch!
\zl
Wird eine andere Konstituente vorangestellt, ergibt sich eine Frage mit Fragewort (\mex{1}a), ein Imperativ
(\mex{1}b) oder ein deklarativer Satz (\mex{1}c).
\eal
\ex 
Wem gibt er das Buch?
\ex 
Jetzt gib ihm das Buch!
\ex 
Jetzt gibt er ihm das Buch.
\zl
Die Analyse der Semantik von Satztypen kann hier nicht gegeben werden, aber die interessierte Leserin oder der interessierte Leser
sei auf \citew{MuellerSatztypen,MuellerGS} verwiesen.%
\is{Vorfeldbesetzung|)}

%\if0
\section{Alternativen}

Wie bei den Abschnitten über Alternativen in den vorigen Kapiteln ist dieser Abschnitt nur für fortgeschrittene
Leserinnen und Leser. Es ist nicht nötig, ihn zu lesen, um den Rest des Buches zu verstehen.

Im vorigen Abschnitt habe ich eine Analyse vorgeschlagen, in der die grundlegende SVO-Abfolge genau das ist: ein
Subjekt, gefolgt vom Verb, und ein Verb, gefolgt von den Objekten. Die verb"=finalen Sätze von SOV-Sprachen
werden als ein Verb analysiert, dem seine Argumente vorangehen. Die Position des finiten Verbs
wird dadurch erfasst, dass es über den Doppelschrägstrich-Mechanismus vorangestellt wird.

%\largerpage[1]
Es gibt alternative Vorschläge zur SVO- und SOV-Abfolge und auch zur Platzierung des finiten
Verbs. Der Vorschlag von \citet{Kayne94a-u} besagt, dass alle Sprachen eine zugrundeliegende
Spezifikator-Kopf-Komplement-Abfolge haben. Die Abfolgen, die wir in den \ili{germanisch}en SOV-Sprachen sehen, würden dann
durch Bewegung abgeleitet. Der Gegenvorschlag von \citet{Haider2000a,Haider2020a} besagt nicht, dass alle Sprachen
wie das \ili{Englisch}e oder das \ili{Romanisch}e sind, sondern behauptet stattdessen, dass die VOO-Muster in SVO-Sprachen
aus einem zugrundeliegenden OVO-Muster abgeleitet sind. Diese beiden Ansätze werden in den folgenden zwei Unterabschnitten~\ref{sec-ov-derived-from-vo} und~\ref{sec-vo-derived-from-ov} diskutiert. Wie gezeigt wird,
macht Kaynes Vorschlag
%\pagebreak{}
falsche Vorhersagen, und Haiders Vorschlag ist ebenfalls nicht ohne Probleme. Bei
beiden Vorschlägen wäre unklar, wie sie von Lernenden der jeweiligen
Sprachen ohne die Annahme einer reichhaltigen Universalgrammatik erworben werden sollten.

Die dritte Klasse von Vorschlägen, die in Abschnitt~\ref{sec-aux-flat} diskutiert wird, nimmt überhaupt keine
Verbbewegung an. Statt eine Struktur mit geschichteten VPs und einer Art Bewegung anzunehmen, die das
finite Verb umstellt, nehmen Autoren wie \citet*{GKPS85a} und \citet{Sag2020a} an, dass es alternative
Linearisierungen für finite Verben und ihre Subjekte gibt. Die Vor- und Nachteile solcher Analysen sind das Thema
von Abschnitt~\ref{sec-aux-flat}.

Das CP/TP/VP-Modell wurde bereits in Abschnitt~\ref{sec-cp-tp-vp-scrambling} über Scrambling diskutiert.
Es gibt auch Argumente gegen diesen Ansatz, was die Verbbewegung betrifft. Sie werden in
Abschnitt~\ref{sec-cp-tp-vp-and-verb-last-non-movement} diskutiert.

\subsection{OV abgeleitet aus VO: \citet{Kayne94a-u}}
\label{sec-ov-derived-from-vo}

%\largerpage
\citet{Kayne94a-u} postuliert, dass alle Sätze in allen Sprachen eine Spezifikator–Kopf–Komplement-Abfolge
haben. Bei Sprachen mit Abfolgen, die nicht SVO sind, wird angenommen, dass sie durch Bewegung aus SVO abgeleitet sind. Wir
haben \citegen[\page 224]{Laenzlinger2004a} Analyse eines \ili{deutsch}en Satzes bereits in
Abschnitt~\ref{sec-autonomy-of-syntax} diskutiert. Abbildung~\ref{fig-Kayne-for-German} zeigt seine Analyse ohne
die funktionalen Knoten für Adjunkte und ohne die Voranstellung des Objekts zu TopP.

\begin{figure}
%\oneline{%
\scalebox{.93}{%
\begin{forest}
%where n children=0{delay=with unaligned translation}{}
sm edges
[CP
	[C$^0$[weil, tier=word]]
%	\[TopP
%		[DP$_j$ [diese Sonate;this sonata]]
		[SubjP
			[DP$_i$ [der Mann,roof]]
%			\[ModP
%				[AdvP [wahrscheinlich;probably,l=20\baselineskip]]
				[ObjP
					[DP$_j$ [diese Sonate,roof]]
%					\[NegP
%						[AdvP [nicht;not,tier=word]]
%						\[AspP
%							[AdvP [oft;often,tier=word]]
%							\[MannP
%								[AdvP [gut;well,tier=word]]
								[AuxP
									[VP$_k$ [gespielt,tier=word]]
									[Aux+
										[Aux [hat,tier=word]]
										[vP
											[DP$_i$ [,phantom]]
											[VP$_k$
												[V [,phantom]]
												[DP$_j$
                                                                                                  [,phantom]]]]]]]]]
%]]]]
\end{forest}%
}
\caption{\label{fig-Kayne-for-German}Verkürzte Analyse der Satzstruktur mit Restbewegung nach links
  und funktionalen Köpfen nach \citet[\page 224]{Laenzlinger2004a}}
\end{figure}%

Die Abbildung zeigt eine vP rechts vom Auxiliar. Es wird also angenommen, dass die zugrundeliegende Struktur ohne Adjunkte
(\mex{1}a) ist und die abgeleitete (\mex{1}b):
\eal
\ex[*]{ 
weil hat der Mann gespielt die Sonate
}
\ex[]{
weil der Mann die Sonate gespielt hat
}
\zl

\noindent
Es gibt ein sehr einfaches Argument gegen solche Ableitungen: Es stammt aus dem Poverty-of-the-Stimulus-Argument,
einem alten Argument von \citet[\page 34]{Chomsky80b-u}, und ich möchte das Argument, das ich
hier verwende, das inverse Poverty-of-the-Stimulus-Argument\is{Poverty of the stimulus}\label{pos-argument} nennen, da ich dasselbe Argument in eine
andere Richtung verwende. Chomsky argumentiert, dass Wissen, das nicht aus dem
Input erlernt werden kann und dennoch nachweislich vorhanden ist, angeboren sein muss. Da wir inzwischen wissen, dass es
höchst unwahrscheinlich ist, dass ausgearbeitetes sprachspezifisches Wissen Teil unseres Genoms ist
% due to a page break this would crash lualatex as of 7/2025
(\citealt*{HCF2002a}, \cites{Bishop2002a}[Section~6.4.2.2]{Dabrowska2004a}, \citeauthor{FM2005a} \citeyear{FM2005a}), folgt daraus, dass
die in der Linguistik angenommene Maschinerie nicht so beschaffen sein kann, dass sie nicht erwerbbar wäre. Da die
von \citeauthor{Laenzlinger2004a} angenommene zugrundeliegende Struktur in keiner Weise mit beobachtbarem
Material verbunden ist, gäbe es keine Möglichkeit, die Transformationen zur Ableitung \ili{deutsch}er
Sätze zu erlernen. Daraus folgt, dass die vollständige Maschinerie Teil der Universalgrammatik sein müsste, aber
da unklar ist, wie sie dorthin gelangen sollte und warum man sie überhaupt annehmen sollte,
müssen wir schlussfolgern, dass \citeauthor{Kayne94a-u}s Vorschlag falsch ist.

%\largerpage[1.1]
\citet{Haider2000a} zeigt im Detail, warum man nicht annehmen sollte, dass OV aus VO abgeleitet ist. Ich werde die
Diskussion hier nicht wiederholen. Haider argumentiert, dass man stattdessen VO als aus OV abgeleitet ansehen sollte. Ich denke,
dass auch dies keine gute Idee ist und dass VO und OV einfach verschieden und nicht voneinander
abgeleitet sind. Ich befasse mich mit Haiders Vorschlag und meinem alternativen Vorschlag im nächsten Unterabschnitt.

\subsection{VOO abgeleitet aus OVO: \citet{Haider2020a}}
\label{sec-vo-derived-from-ov}

Wie im vorigen Abschnitt gezeigt wurde, sind SVO-Ansätze für SOV-Sprachen nicht auf
oberflächenorientierte Weise erwerbbar, erfordern einen großen Teil angeborenen sprachspezifischen Wissens und sind daher
mit allem, was wir über den Spracherwerb wissen, unvereinbar. Nun argumentiert Hubert
\citet{Haider2000a,Haider2010a,Haider2020a} in Bezug auf die Psycholinguistik, dass eine VO-Sprache wie das
\ili{Englisch}e eine Struktur hat, in der das Verb mit Argumenten auf eine ähnliche Weise kombiniert wird wie in OV-Sprachen
wie dem \ili{Deutsch}en, außer dass der Kopf in den VO-Sprachen in eine Anfangsposition bewegt wird.
\citet[\page 15]{Haider2010a} vergleicht kopfinitiale und kopffinale Ansätze:
\eal
\ex {}[[[h$^0$ A$_1$] A$_2$] A$_3$]
\ex {}[A$_3$ [A$_2$ [A$_1$ h$^0$]]]
\zl
Er argumentiert, dass (\mex{0}a), das das Spiegelbild von (\mex{0}b) ist, sprachübergreifend nicht existiert. Er argumentiert, dass die Argumentstruktur
in VO- und OV-Sprachen dieselbe ist und dass VO-Sprachen dieselbe Abfolge der Argumente haben wie die OV-Sprachen.
Er schließt auf S.\,28, dass die Satzstruktur mit einem \ili{englisch}en dreistelligen Verb wie in (\mex{1}) aussieht:
\ea
\label{ex-Haider-English-head-movmenet}
{}[XP [h$^0$ [YP [h$^0$ ZP]]]]
\z
Das bedeutet, dass er annimmt, dass (\mex{1}a) aus dem zugrundeliegenden (\mex{1}b) abgeleitet ist:
\eal
\ex[]{
Kim gave Sandy a book.
}
\ex[*]{
Kim Sandy gave a book.
}
\zl
Wieder greift das inverse Poverty-of-the-Stimulus-Argument\is{Poverty of the stimulus} (siehe S.\,\pageref{pos-argument}): Die Lernerin oder der Lerner des
Englischen hat keine Möglichkeit herauszufinden, dass (\mex{0}a) mit (\mex{0}b) in Beziehung steht, daher muss dieses Wissen
irgendwie Teil einer angeborenen Komponente der Sprache sein, etwas, das Haider tatsächlich annimmt. Wenn wir
solche komplexe domänenspezifische Information als Teil unserer genetischen Ausstattung ablehnen, müssen wir Haiders Analyse ablehnen.

Ich habe in den vorigen Kapiteln stattdessen argumentiert, dass die Reihenfolge, in der die Argumente mit ihren Verben kombiniert werden,
frei ist. Daher können wir das Verb zuerst mit A$_3$ kombinieren, selbst wenn das Verb kopfinitial ist:
\ea
{}[[[h$^0$ A$_3$] A$_2$] A$_1$]
\z
Für das \ili{Englisch}e erhalten wir eine Satzstruktur wie in (\mex{1}):
\ea
{}[XP [[h$^0$ YP] ZP]]
\z
%\largerpage
Der Unterschied zwischen Haiders Struktur in (\ref{ex-Haider-English-head-movmenet}) und (\mex{0})
ist, dass Haider annimmt, dass es in einfachen \ili{englisch}en SVO-Strukturen Kopf"=bewegung gibt. Der Kopf beginnt links von ZP,
wo er die ZP rechts selegiert, und bewegt sich dann nach oben an die Spitze von YP. Dies wird in der linken
Abbildung in Abbildung~\ref{fig-Haider-English-German} gezeigt.
\begin{figure}
\hfill
\begin{forest}
[vP
  [XP]
  [v$'_{\nliste{x}}$
     [v$_i^0$,name=v2]
     [VP$_{\nliste{x}}$
       [YP,name=YP]
       [V$'_{\nliste{x,y}}$
          [V$_{i\nliste{x,y,z}}^0$,name=v1]
          [ZP]]]]]
\draw[->](v2) to (YP);
\draw[->](YP) to (v1);
\end{forest}
\hfill
\begin{forest}
[VP
  [XP,name=XP]
  [V$'_{\nliste{x}}$,name=V2
     [YP,name=YP]
     [V$'_{\nliste{x,y}}$,name=V1
        [ZP,name=ZP]
        [V$_{i\nliste{x,y,z}}^0$,name=V0]]]]
% \draw[->](V0.west) to (ZP.east);
% \draw[->](V1.west) to (YP.east);
% \draw[->](V2.west) to (XP.east);
\draw[->](V0) to (ZP);
\draw[->](V1) to (YP);
\draw[->](V2) to (XP);
\end{forest}
\hfill\mbox{}
\caption{Englisch vs. Deutsch nach \citet[\page 29]{Haider2010a}}\label{fig-Haider-English-German}

\end{figure}

Haider argumentiert, dass die Strukturen für das \ili{Englisch}e durch die UG bestimmt werden und dass es zu viel
Struktur mit zu vielen Klammern gäbe, wenn es keine Kopf"=bewegung gäbe, um die Struktur
aus einer Verarbeitungsperspektive plausibel zu machen.

Die Frage ist: Was ist eine psycholinguistisch plausible Geschichte für die Analyse \ili{englisch}er
Sätze? Wir wissen jetzt mit Sicherheit, dass die menschliche Sprachverarbeitung inkrementell ist \citep{Marslen-Wilson75a,TSKES96a,SW2011a,Wasow2021a}. Wir verwenden Information
aus allen verfügbaren Quellen, sobald wir sie haben: Phonologie, Syntax, Semantik, Pragmatik,
Gesten, Weltwissen. Wenn wir eine NP hören, stellen wir Hypothesen darüber an, wie
die Äußerung weitergehen könnte. Eine mögliche Fortsetzung ist die als Satz. Wir erwarten also eine VP, die der
NP folgt, wie in Abbildung~\ref{fig-NP-S}.\footnote{
Siehe auch \citet{Jurafsky96a} für eine Erklärung von Holzwegsätzen
mithilfe eines probabilistischen Chart-Parsers und einer Phrasenstrukturgrammatik mit Merkmal-Wert-Paaren der Art,
wie sie in HPSG angenommen werden. Jurafsky verbindet Wahrscheinlichkeiten mit Phrasenstrukturregeln und mit Valenz\-information.
Für die hier entwickelte Grammatik spielen Valenzrahmen eine entscheidende Rolle bei der Vorhersage
kommender Konstituenten.}
\begin{figure}
\begin{forest}
[S
  [NP
    [Kim,tier=word]]
  [VP,edge=dashed
    [V,edge=dashed [,no edge,tier=word]]]]
\end{forest}
\caption{Es wird erwartet, dass eine VP der NP folgt, um einen Satz zu bilden}\label{fig-NP-S}
\end{figure}
Nachdem wir eine NP gehört haben, erwarten wir eine VP, die irgendwo ein Verb enthält, aber das nächste kommende Wort könnte
auch ein Adverb oder zwei sein, wie in (\mex{1}a):
\eal
\ex their willingness [[usually [strongly [depends on this]]]]\footnote{
  ENCOW, doc\#40288,www.psy.gla.ac.uk
}
\ex Kim [[promised and gave] Robin a book].
\zl
%%\largerpage
Es könnte auch ein Verb sein, das Teil einer Koordination von zwei oder mehr Verben ist. Und Kombinationen sind
natürlich möglich. Die Tatsache, dass die Struktur unterdeterminiert ist, wird durch eine gestrichelte Linie angezeigt.
Wenn wir das nächste Wort hören, können wir allerdings eine konkretere Hypothese bilden. Im
Falle eines strikt transitiven Verbs könnten wir es als Teil einer Koordination haben oder, wahrscheinlicher, als Teil
einer VP, die aus einem Verb und einer NP besteht. Die nächste NP wird vorhergesagt wie in Abbildung~\ref{fig-NP-S-V}.
\begin{figure}
\begin{forest}
[S
  [NP
    [Kim,tier=word]]
  [VP,edge=dashed
    [V
      [devoured,tier=word]]
    [NP]]]
\end{forest}
\caption{Es wird erwartet, dass eine NP der NP + Verb folgt, um mit einem strikt transitiven Verb einen Satz zu bilden}\label{fig-NP-S-V}
\end{figure}
Wenn wir statt eines strikt transitiven Verbs ein ditransitives Verb hören, wird stattdessen die Struktur in Abbildung~\ref{fig-NP-S-V-ditrans}
vorhergesagt.\footnote{
  Man beachte, dass es aus psycholinguistischer Sicht keinen Unterschied zwischen dieser binär
  verzweigenden VP-Struktur und einer flachen gibt. In einem Ansatz mit flacher Struktur würde man annehmen, dass man
  eine VP begonnen hat und in einer flachen Struktur eine Tochter nach der anderen prüfen würde, während man in einem binär verzweigenden Ansatz
  ausgefeiltere Strukturen hat und die NP-Töchter in separaten Teilbäumen
  nacheinander prüft.
}
\begin{figure}[t]
\begin{forest}
[S
  [NP
    [Kim,tier=word]]
  [VP,edge=dashed
    [V$'$
      [V
        [gave,tier=word]]
      [NP]]
     [NP]]]
\end{forest}
\caption{Es wird erwartet, dass zwei NPs der NP + Verb folgen, um mit einem ditransitiven Verb einen Satz zu bilden}\label{fig-NP-S-V-ditrans}
\end{figure}
Die Objekt-NPs müssen eingefügt werden wie in Abbildung~\ref{fig-NP-S-V-ditrans-NP-NP}, aber weitere Adjunkte können rechts von der VP hinzugefügt werden. Es
muss also Platz dafür sein. Alles, was wir bis zum Ende des Satzes mit Sicherheit wissen, ist, dass der S-Knoten
eine VP dominiert. Es kann mehr als einen VP-Knoten geben.
\enlargethispage{3pt}
\begin{figure}[htb]
\begin{forest}
sm edges
[S
  [NP
    [Kim]]
  [VP,edge=dashed
    [V$'$
      [V
        [gave]]
      [NP [Robin]]]
     [NP [the book,roof]]]]
\end{forest}
\caption{NP + ditransitives Verb + zwei Objekte bilden einen Satz, zu dem Adjunkte hinzugefügt werden können.}\label{fig-NP-S-V-ditrans-NP-NP}
\end{figure} 

\begin{figure}[htb]
\begin{forest}
sm edges
[S
  [NP
    [Kim]]
  [VP
    [VP
      [V$'$
        [V
          [gave]]
        [NP [Robin]]]
      [NP [the book,roof]]]
    [Adv [yesterday]]]]
\end{forest}
\caption{Satz mit VP-Adjunkt rechts.}\label{fig-NP-S-V-ditrans-adverb}
\end{figure}

Was all dies uns zeigt, ist, dass Sätze intern komplex sein können und mehrere Adjunkte an eine VP
angeschlossen werden können. Dennoch kann der menschliche Satzverarbeiter damit und mit der standardmäßig
angenommenen Schachtelung von Adverbialen und VPs umgehen. Daraus folgt, dass es unnötig ist, einen
komplizierten Kopf"=bewegungsansatz für \ili{englisch}e verbale Projektionen anzunehmen.

%\largerpage
Shravan Vasishth erinnert mich daran, dass „absence of evidence is not evidence of absence“. Das bedeutet, dass
die Tatsache, dass wir keine psycholinguistischen Reflexe der von Haider angenommenen Strukturen finden können,
nicht bedeutet, dass sie nicht vorhanden sind. Das stimmt, aber ich möchte argumentieren, dass evidence of absence
zusammen mit Ockhams Rasiermesser ein Argument gegen bestimmte Strukturen ist: Wenn eine Struktur komplex und
unnötig ist und es keine psycholinguistische Evidenz \emph{für} sie gibt, sollte sie nicht angenommen werden.

% \subsection{Binding and branching}

% The previous section discussed \citeauthor{Haider2010a}'s approach, which assumes that VO languages
% are derived from OV languages. The structures he assumes for sentences with ditransitive verbs are
% similar to structures assumed elsewhere in GB/""Minimalism. For example, \citet{Adger2003a} assumes
% the right-most structure in Figure~\ref{fig-ditransitives-options}.
% \begin{figure}
% \begin{forest}
% baseline
% [\vbar
%  [\textit{show}]
%  [\textit{himself}]
%  [\textit{Benjamin}]]
% \end{forest}
% \hfill
% \begin{forest}
% baseline
% [\vbar
%    [\vbar
%      [\textit{show}]
%      [\textit{himself}] ]
%  [\textit{Benjamin}]]
% \end{forest}
% \hfill\hfill
% \begin{forest}
% baseline
% [\littlevbar
%  [\textit{show}]
%  [VP
%    [\textit{himself}]
%    [\vbar
%     [V]
%     [\textit{Benjamin}]]]]
% \end{forest}
% \caption{\label{fig-ditransitives-options}Three possible analyses of ditransitives}
% \end{figure}
% The motivation for this structure is Binding Theory. Binding Theory deals, among other things, with
% the distribution of personal pronouns and reflexives. It was an important part of Chomsky's book on
% Government \& Binding \citep{Chomsky81a}. \ili{English} reflexives have to match their antecedent in
% gender:
% \eal
% \ex Peter saw himself.
% \ex Mary saw herself.
% \zl
% Reflexives must be bound in a certain local domain, while personal pronouns have to be free. For
% example, the pronouns in (\mex{1}) cannot refer to Peter and Mary, but they may refer to John and
% Helen, respectively.
% \eal
% \ex John claims that Peter saw him.
% \ex Helen claims that Mary saw her.
% \zl
% For the definition of binding domains the notion of c-command plays an important role. c-command is
% defined with respect to trees. (\mex{1}) gives the definition that is used by \citet[\page 117]{Adger2003a}:
% \ea
% A node A c-commands B if, and only if A's sister either:\\
% \begin{tabular}[t]{@{}l@{~}l@{}}
% a. & is B, or\\
% b. & contains B
% \end{tabular}
% \z

% The following two sentences show that \emph{himself} cannot be bound by Benjamin. The first example
% is simply ungrammatical since \emph{Emily} is not a suited candidate for binding because of its
% gender and the second example is grammatical but only with a binding of \emph{himself} by Peter.
% \eal
% \ex[*]{
% Emily showed himself Benjamin in the mirror.
% }
% \ex[]{
% \label{ex-Peter-showed-himself-Benjamin}
% Peter showed himself Benjamin in the mirror.
% }
% \zl
% In order to be able to explain these binding facts via c-command, a structure like the left one or
% the one in the middle of Figure~\ref{fig-ditransitives-options} would not be
% suited. \citeauthor{Adger2003a} uses a special functional head \littlev and the 
% verb \emph{show} moves from V to \littlev \citep{Larson88a}. \Littlev is assumed to contribute a causative meaning
% component \citep[\page 70]{HK93a-u}.

% The theory presented here assumes the structure in the middle. It is a binary branching structure,
% but there is no head"=movement in simple SVO structures. The verb is just combined with one object
% after the other.\footnote{
%   This kind of branching is also assumed in Categorial Grammar. See for example \citew[\page 392]{Steedman2019a-u}.
% }
% Now, there seems to be something missing. How is binding accounted for? c-command cannot be used for
% this. There is an interesting solution that does not require trees that are structured in the
% GB/""Minimalism way: \textcites{PS92a,ps2} developed a Binding Theory that is based on the order of
% elements in the \argstl. The \argstl basically also determines local trees containing the subject
% and objects. Since the elements in the \argstl are ordered in a certain way, one can use this list
% for establishing or ruling out binding relations. Trees are not needed for binding. This means that
% HPSG trees can be formed according to the adjacent material that is actually combined. Movement and
% additional structure is not necessary. For further details on Binding Theory in HPSG and a more
% recent approach see \citew{Branco2002a}. An overview of HPSG's Binding Theory can be found in
% \citew{Mueller2021a}.

\subsection{Analysen verb-initialer Sätze in SVO-Sprachen ohne Verbbewegung}
\label{sec-aux-flat}

Ich habe in Abschnitt~\ref{sec-phenomenon-v2} erwähnt, dass das \ili{Englisch}e eine residuale V2-Sprache
ist. \emph{wh}-Fragen wie die in (\ref{ex-which-book-did-sandy-give-to-kim}) auf
S.\,\pageref{ex-which-book-did-sandy-give-to-kim}, hier als (\mex{1}) wiederholt, sind in ihrer Form ähnlich zu dem, was wir
für die anderen \ili{germanisch}en Sprachen gesehen haben: Eine \emph{wh}-Phrase wird vorangestellt, und das Subjekt steht in der
Position rechts vom finiten Verb.
\ea
Which book did Sandy give to Kim?
\z
\noindent
\citet{Borsley89} hat eine Verbbewegungsanalyse der \ili{englisch}en Satzstruktur
entwickelt. Abbildung~\ref{fig-verb-movement-English} zeigt die Analyse in der hier verwendeten Notation.

\begin{figure}
\centering
\begin{forest}
sm edges
[S
  [NP$_i$ [which book,roof]]
  [S/NP
    [V \sliste{ S//V }
       [V [did$_j$] ] ]
          [S//\vSLASHnp
             [NP [Sandy]]
             [VP//\vSLASHnp
               [V//V  [\_$_j$] ]
               [VP/NP, s sep*=2
                 [\vbarSLASH NP
                   [V [give]]
                   [NP/NP [\trace$_i$ ]]]
                 [PP [to Kim,roof]]]]]]]
\end{forest}
\caption{\label{fig-verb-movement-English}Analyse von \emph{wh}-Interrogativsätzen nach \citet{Borsley89}}
\end{figure}%

Obwohl diese Analyse mit dem kompatibel ist, was an anderer Stelle in diesem Buch gesagt wird, ist sie nicht die Analyse,
die üblicherweise in HPSG angenommen wird. Die Standardanalyse geht auf \citet*{GKPS85a} zurück und wird von
mehreren Autoren in verschiedenen Formen aufgegriffen. Die jüngste Analyse wurde von Ivan Sag entwickelt und posthum
als \citew{Sag2020a} veröffentlicht. Sag schlägt eine flache Analyse \ili{englisch}er VPs vor. Das bedeutet, dass das
Verb und alle Objekte Töchter derselben Mutter sind. Für Fälle von Subjekt-Verb-Inversion gibt es
mehrere Vorschläge. Abbildung~\ref{fig-did-kim-get-the-job-Sag} zeigt eine flache Analyse \ili{englisch}er
verb-initialer Fragen: Sowohl das Subjekt als auch das VP-Komplement sind Mitglieder der \compsl und können
mit dem Verb in Richtung der Komplemente kombiniert werden.
\begin{figure}
\begin{forest}
sm edges
[S
  [{V[\comps \sliste{ \ibox{1}, \ibox{2} } ]} [did]]
  [\ibox{1} NP [Kim]]
  [\ibox{2} VP [get the job, roof]]]
\end{forest}
\caption{\label{fig-did-kim-get-the-job-Sag}Analyse englischer Auxiliarkonstruktionen auf der Basis von
  \citet[\page 117]{Sag2020a}}
\end{figure}

Bei einem solchen Ansatz ist Verbbewegung unnötig. Der Vorteil solcher flachen Ansätze ist, dass
verschiedene spezialisierte Bedeutungen von Auxiliarinversionskonstruktionen dieser
Konfiguration zugewiesen werden können. Zum Beispiel postuliert \citet[\page 116]{Sag2020a} einen Subtyp \type{polar-int-cl} für polare
Interrogativsätze wie (\mex{1}a) und einen weiteren Subtyp \type{aux-initial-excl-cl} für Exklamativsätze wie (\mex{1}b).
\eal
\ex Are they crazy?
\ex Are they crazy!
\zl

\noindent
In dem hier entwickelten Ansatz muss die jeweilige Bedeutung den Auxiliaren zugewiesen werden.

Es ist schwierig, zwischen den beiden Ansätzen zu entscheiden: Einerseits ist die flache Analyse einfacher
als die mit Verbbewegung. Da es keine Evidenz für eine reichhaltige Universalgrammatik im
Chomskyschen Sinne gibt \citep{HCF2002a}, kann man nicht annehmen, dass das \ili{Englisch}e dieselbe Struktur hat wie die
\ili{germanisch}en V2-Sprachen, da wir uns nicht darauf verlassen können, dass Wissen über diese Strukturen
angeboren ist. Sprachen müssen auf der Grundlage von Daten der jeweils betrachteten Sprache
allein analysiert werden, da dies die Daten sind, die Sprachlernenden zur Verfügung stehen. Aus dieser Perspektive gewinnt die flache
Analyse. Andererseits sollte man Analysen anstreben, die über Sprachen hinweg ähnlich sind. Dies
wäre ein Argument für die Verbbewegungsanalyse. Ich habe in \citew{MuellerCoreGram} und in
\citew{MuellerGT-Eng5} argumentiert, dass Grammatiken auf der Grundlage von Daten der jeweils betrachteten Sprache
allein geschrieben werden sollten. Wenn wir mehrere Optionen haben, ein bestimmtes Phänomen zu analysieren, können wir die
Analyse wählen, die für mehrere Sprachen funktioniert, und so sprachübergreifende Generalisierungen erfassen. Im
Fall der englischen Auxiliarinversion scheint es einen sprachinternen Vorteil der
Nicht-Bewegungs-Analyse zu geben. Es bleibt abzuwarten, ob es Daten gibt, die uns zwingen, eine Verbbewegungsanalyse
anzunehmen, wie es bei den Daten zur multiplen Voranstellung im \ili{Deutsch}en der Fall war \citep{Mueller2005d,MuellerGS}.

\subsection{CP/TP/VP}
\label{sec-cp-tp-vp-and-verb-last-non-movement}

Abschnitt~\ref{sec-cp-tp-vp} war der CP/TP/VP-Analyse und dem
Scrambling gewidmet. Hier diskutiere ich Doppelpartikelverben (Abschnitt~\ref{sec-cp-tp-vp-double-particles}) und Landeplätze für die Extraposition (Abschnitt~\ref{sec-cp-tp-vp-extraposition}).

\subsubsection{Doppelpartikelverben}
\label{sec-cp-tp-vp-double-particles}

% Haider 1993: 62; Vikner 2001; Fortmann 2007)

%\largerpage
Abbildung~\ref{fig-cp-tp-vp} auf S.\,\pageref{fig-cp-tp-vp} zeigte, dass angenommen wird, dass sich der Verbstamm von V nach T bewegt, um
dort eine Endung aufzunehmen und Flexionsmerkmale zu überprüfen. \textcites[\page 63]{Haider93a}[\page
59--60]{Haider2010a} und \citet[Section~3.3]{Vikner2001a} haben ein
Argument gegen solche Vorschläge gefunden: Das \ili{Deutsch}e hat bestimmte Verben mit zwei Partikeln.\footnote{%
  Die Beobachtung, dass das \ili{Deutsch}e Verben hat, die nicht an V2-Sätzen teilnehmen können, geht auf
  \citet{Hoehle91b} zurück.
}
\emph{vorankündigen}
ist ein Beispiel. Dieses Partikelverb besteht aus der Kombination von \emph{an} und \emph{kündigen} mit
der zusätzlichen Hinzufügung eines weiteren Präfixes \emph{vor}. Beispiel
(\ref{ex-dass-sie-es-vorankündigt}) wurde bereits auf
S.\,\pageref{ex-dass-sie-es-vorankündigt} diskutiert, aber es wird hier der Bequemlichkeit halber als (\mex{1}) wiederholt:
\ea 
\label{ex-dass-sie-es-vorankündigt-two}
dass sie es vor-an-kündigt
\z
Das Interessante an diesen Verben ist nun, dass sie nicht vorangestellt werden können. Der Verbstamm muss
benachbart zur Partikel stehen:
\eal
\ex[*]{
Sie kündigt$_i$ es vor-an \_$_i$.
}
\ex[*]{
Sie an-kündigt$_i$ es vor \_$_i$.
}
\ex[*]{
Sie vor-an-kündigt$_i$ es  \_$_i$.
}
\zl
Haider hat darauf hingewiesen, dass für solche Verben unter Ansätzen, die annehmen, dass sich der Verbstamm von V nach T bewegt, um
\isi{Kongruenz}merkmale zu überprüfen, vorhergesagt wird, dass sie keine finiten Formen haben.

Die TP-basierte Analyse müsste annehmen, dass sich der Verbstamm \stem{kündig} von V nach T bewegt wie in (\mex{1}), aber
da diese Art der Bewegung für Doppelpartikelverben ausgeschlossen ist, wie (\mex{0}) zeigt, sollten finite Formen
von Doppelpartikelverben nicht existieren, nicht einmal in satzfinaler Position.
\ea 
dass sie es vor-an \_$_i$ kündig$_i$ -t
\z
Aber wie (\ref{ex-dass-sie-es-vorankündigt-two}) zeigt, existieren solche Sätze. Wie Haider also
hervorgehoben hat, scheint ein CP/VP-Modell angemessener zu sein. Verben bewegen sich nicht in höhere funktionale Projektionen
wie T, um ihre Kongruenzmerkmale zu überprüfen. Sie tun das einfach in der Position, in der sie sind: in der V-Position. Die
Annahme einer T-Projektion für \isi{Kongruenz} ist unnötig, ja sie ist mit den
beobachtbaren Daten unvereinbar.\footnote{%
  \citet{Vikner2001a} nimmt eine TP für alle \ili{germanisch}en Sprachen an, sogar für OV-Sprachen. Er nimmt
  an, dass es in diesen Sprachen keine V-zu-T-Bewegung gibt. Modelle, die diese
  zusätzliche Schicht einfach nicht annehmen, scheinen jedoch angemessener zu sein.
}

\subsubsection{VPs als Landeplätze}
\label{sec-cp-tp-vp-extraposition}

% Hubert Haider noted early on that sentences like (\mex{1}) pose a problem for transformational
% theories:
% \eal
% \ex[]{
% Einen Hund füttern, der Hunger hat, wird man müssen.
% }
% \ex[*]{
% Wird man einen Hund füttern, der Hunger hat, müssen?
% }
% \ex[]{
% Wird man einen Hund füttern müssen, der Hunger hat?
% }
% \zl

%\largerpage
\citet[\page 62--64]{Haider2010a} hat Daten zur Platzierung von PPs/Adverbien in Bezug auf die V-zu-T-Bewegungshypothese
untersucht. Er hat festgestellt, dass präpositionale Adverbiale nur marginal zwischen die Verben eines Verbalkomplexes
platziert werden können und dass vollständige PPs ungrammatisch sind. Sowohl PPs als auch Pronominaladverbien sind zwischen
Auxiliaren und Modalverben vollständig inakzeptabel:
\eal
\ex[?/*]{
dass er viel gelernt \emph{dafür} haben muss
}
\ex[*]{
dass er viel gelernt haben \emph{dafür} muss
}
\ex[?/*]{
ohne    viel gelernt \emph{dafür} haben zu müssen
}
\ex[*]{
ohne    viel gelernt haben \emph{dafür} zu müssen
}
\ex[*]{
dass er viel gelernt \emph{für} \emph{das} \emph{Examen} hat
}
\ex[*]{
ohne    viel gelernt \emph{für} \emph{das} \emph{Examen} zu haben
}
\ex[]{
[\sub{VP} Gelernt haben \emph{dafür} / \emph{für} \emph{das} \emph{Examen}] muss er viel.
}
\zl 
Wie (\mex{0}g) jedoch zeigt, sind VPs ein legitimer Landeplatz für die PP-Extraposition und für die
  Extraposition von Pronominaladverbien. \emph{Gelernt haben} bildet die rechte Satzklammer,
  und das PP-Material wird im \nf der vorangestellten VP platziert.
Haider fährt mit den Beispielen in (\mex{1}) fort:
\eal
\ex[]{
[\sub{VP} Angefangen damit]$_i$ hat bloß einer \trace$_i$
}
\ex[*]{
weil    bloß einer an \trace$_i$ damit   fing$_i$
}
\ex[]{
weil bloß einer anfing damit
}
\zl
(\mex{0}a) zeigt, dass \emph{damit} rechts von einem Partikelverb platziert werden kann. Wenn es eine VP gibt,
die unter einer TP eingebettet ist, würde man erwarten, dass diese VP ebenfalls ein möglicher Landeplatz
für die Extraposition ist, wie es in (\mex{0}a) der Fall ist. Man würde erwarten, dass \emph{damit} neben
der Verbpartikel \emph{an} platziert werden kann wie in (\mex{0}b), aber wenn das Pronominaladverb extraponiert wird,
muss es rechts vom Verb stehen wie in (\mex{0}c).\footnote{
Man beachte, dass die Situation nicht so einfach ist, wie man erwarten könnte. Partikeln können im
\mf des \ili{Deutsch}en platziert werden. Aber diese Partikelplatzierungen unterscheiden sich von der Bewegung des Verbs nach
rechts. Solche Bewegungen müssen in jedem Fall ausgeschlossen werden.%
} Es mag Wege geben, die problematischen Daten in einem
VP/TP-System wegzuerklären, aber die einfachste Erklärung ist natürlich, von vornherein keine TP
anzunehmen. Wenn es ein Verb \emph{anfing} als Kopf der VP gibt, bewegt es sich in
verbletzten Sätzen nicht zu einem höheren Kopf, und daher muss kein Material daran gehindert werden, zwischen VP und T zu intervenieren.

\clearpage

\questions{

\begin{enumerate}
\item Wie werden Satztypen in den \ili{germanisch}en Sprachen bestimmt?

\end{enumerate}

}

\exercises{

\begin{enumerate}
\item Klassifizieren Sie die \ili{germanisch}en Sprachen nach ihrer Grundkonstituentenabfolge (SVO, SOV, VSO,
  \ldots) und V2 unter der Annahme, dass Sie wissen, dass eines der folgenden Muster in der Sprache existiert:
\eal
\ex NP[acc] V-Aux NP[nom] V NP[dat]  % V2 SVO cannot be \ili{English} since \ili{English} does not have a dative
\ex NP[acc] V-Aux NP[nom] NP[dat] V  % V2 SOV
\ex NP[acc] NP[nom] V NP[acc]        % -V2 SVO
\ex NP[acc] NP[nom] V-Aux V NP[acc]  % -V2 SVO
\ex NP[acc] V-Aux NP[nom] V PP       % kann man nicht sagen
\zl
Jeder Satz sollte mit ±V2 und einer der sechs Permutationen von S, O und V gepaart werden.

Wenn Sie die Abfolge nicht eindeutig bestimmen können, sagen Sie das bitte. Wenn Sie denken, dass dieses Muster
in keiner der \ili{germanisch}en Sprachen existiert, sagen Sie das. Bitte bedenken Sie, dass das \ili{Englisch}e eine sogenannte
residuale V2-Sprache ist, was bedeutet, dass in der Grammatik einige Spuren von V2 verblieben sind. Denken Sie an
die Fragebildung im \ili{Englisch}en.

\item Skizzieren Sie die Analyse für die folgenden Beispiele. Verwenden Sie die in diesem Kapitel benutzten Abkürzungen,
  das heißt, gehen Sie nicht auf die Details bezüglich \spr und \compsvs ein, sondern verwenden Sie S, VP und V$'$. Verbbewegung sollte mit dem Symbol `//' angezeigt werden.
\eal
\ex
\gll Arbejder Bjarne ihærdigt  på bogen?\\
     arbeitet Bjarne ernsthaft an Buch.\textsc{def}\\\danish
\glt `Arbeitet Bjarne ernsthaft an dem Buch?'
\ex
 Arbeitet Bjarne ernsthaft an dem Buch?
\ex
 Wird sie darüber    nachdenken?
\zl

\item Skizzieren Sie die Analyse für die folgenden Beispiele. Verwenden Sie die Valenzmerkmale \spr und \comps
  anstelle der Abkürzungen S, VP und V$'$. Da der Wert von \spr im \ili{Deutsch}en immer die
  leere Liste ist, können Sie ihn in den \ili{deutsch}en Beispielen weglassen. NPs und PPs können als NP bzw.\ PP
  abgekürzt werden. Verbbewegung sollte mit dem Symbol `//' angezeigt werden.
\eal
\ex 
 dass sie darüber nachdenkt
\ex 
 dass sie darüber nachdenken wird
\ex
 Wird sie darüber nachdenken?
\zl

\eal
\ex
\gll Arbejder Bjarne ihærdigt  på bogen?\\
     arbeitet Bjarne ernsthaft an Buch.\textsc{def}\\\danish
\glt `Arbeitet Bjarne ernsthaft an dem Buch?'
\ex
 Arbeitet Bjarne ernsthaft an dem Buch?
\zl

\item Skizzieren Sie die Analyse der folgenden Beispiele. NPs dürfen abgekürzt werden. Valenzmerkmale sollten
  nicht angegeben werden, sondern stattdessen sollten Knotenbezeichnungen wie V, V$'$, VP und S verwendet werden. Wenn Fern\-abhängigkeiten
  beteiligt sind, zeigen Sie sie mit dem Symbol `/' an.
\eal
\ex Such books, I like.
\ex 
 Solche Bücher mag ich.
\ex
\gll Boger som det elsker jeg.\\
     Bücher wie diese liebe ich\\\danish
\glt `Ich mag solche Bücher.'
\zl

\item Skizzieren Sie die Analyse des Beispiels in (\mex{0}b) und beziehen Sie das Merkmal \comps ein.

\end{enumerate}

}

\furtherreading{
Die Analyse von Fernabhängigkeiten wurde von \citet{Gazdar81a} entwickelt. \citet{Sag2010b} behandelt
weitere Beschränkungen, die in einer Theorie der Fernabhängigkeiten im \ili{Englisch}en nötig sind, und wie sie
in HPSG repräsentiert werden können. Das HPSG-Handbuch enthält auch ein Kapitel über Konstituentenabfolge und
Verbbewegung \citep{MuellerOrder} sowie über Fernabhängigkeiten \citep{BC2021a}. \citet{MuellerGS}
behandelt die \ili{deutsch}e Satzstruktur innerhalb von HPSG und diskutiert verschiedene alternative Ansätze, die
hier nicht diskutiert werden konnten (Generalized Phrase Structure Grammar, Dependenzgrammatik, Konstruktionsgrammatik,
linearisierungsbasierte HPSG-Ansätze).

}

%      <!-- Local IspellDict: de_DE -->
